\documentclass[11pt]{article}

% Change "review" to "final" to generate the final (sometimes called camera-ready) version.
% Change to "preprint" to generate a non-anonymous version with page numbers.
\usepackage[final]{acl}

% Standard package includes
\usepackage{times}
\usepackage{latexsym}

% For proper rendering and hyphenation of words containing Latin characters (including in bib files)
\usepackage[T1]{fontenc}
% For Vietnamese characters
% \usepackage[T5]{fontenc}
% See https://www.latex-project.org/help/documentation/encguide.pdf for other character sets

% This assumes your files are encoded as UTF8
\usepackage[utf8]{inputenc}

% This is not strictly necessary, and may be commented out,
% but it will improve the layout of the manuscript,
% and will typically save some space.
\usepackage{microtype}

% This is also not strictly necessary, and may be commented out.
% However, it will improve the aesthetics of text in
% the typewriter font.
\usepackage{inconsolata}

%Including images in your LaTeX document requires adding
%additional package(s)
\usepackage{graphicx}
\usepackage[most]{tcolorbox}

% -------------
\usepackage{kotex} % (폰트/locale 관련은 가장 먼저)
% 수식, 그림, 표 전처
\usepackage{amsmath}
\usepackage{graphicx}
\usepackage{float}
\usepackage{caption}
\usepackage{amssymb}

\usepackage{comment}
\usepackage{hyperref}
\usepackage{booktabs}
\usepackage{tabularx} % 열 너비 자동 조절
\usepackage{array}
\usepackage{longtable}
\usepackage{ragged2e}
\usepackage{makecell}

\newcolumntype{L}[1]{>{\RaggedRight\arraybackslash}p{#1}}
\newcolumntype{Y}{>{\RaggedRight\arraybackslash}X}

\newcolumntype{Y}{>{\RaggedRight\arraybackslash}X}
\newcolumntype{C}[1]{>{\centering\arraybackslash}p{#1}}

\newtcolorbox{promptbox}[1]{
    breakable,
    enhanced jigsaw,
    width=\linewidth,
    colback=white,
    colframe=black,
    coltitle=white,
    colbacktitle=black,
    title=\textbf{#1},
    fonttitle=\bfseries,
    boxrule=0.8pt,
    arc=2mm,
    left=6pt,
    right=6pt,
    top=5pt,
    bottom=5pt,
    before skip=6pt,
    after skip=6pt
}


\definecolor{caseblue}{HTML}{1F77B4}
\definecolor{casered}{HTML}{D62728}
\definecolor{casegray}{HTML}{F7F7F7}

\newcommand{\cmarkblue}{\textcolor{caseblue}{\ding{51}}}
\newcommand{\xmarkred}{\textcolor{casered}{\ding{55}}}
\newcommand{\goldans}[1]{\textbf{#1}}
\newcommand{\correctans}[1]{\textcolor{caseblue}{\textbf{#1}}}
\newcommand{\wrongans}[1]{\textcolor{casered}{\textbf{#1}}}
\newcommand{\caselabel}[1]{\noindent\textbf{\small #1}\vspace{0.25em}}


\newtcolorbox{qualcasebox}[1]{
    breakable,
    enhanced,
    width=\linewidth,
    colback=white,
    colframe=black!45,
    boxrule=0.55pt,
    arc=1.2mm,
    left=7pt,
    right=7pt,
    top=6pt,
    bottom=6pt,
    before skip=7pt,
    after skip=8pt,
    borderline west={2.2pt}{0pt}{black!65},
    title={\textbf{#1}},
    coltitle=black,
    colbacktitle=gray!10,
    fonttitle=\bfseries
}

\newtcolorbox{maincasebox}{
    breakable,
    enhanced,
    width=\linewidth,
    colback=white,
    colframe=black!45,
    boxrule=0.55pt,
    arc=1.2mm,
    left=7pt,
    right=7pt,
    top=6pt,
    bottom=6pt,
    before skip=7pt,
    after skip=8pt,
    borderline west={2.2pt}{0pt}{black!65}
}


% 색/표 스타일 — xcolor는 한 번만!
\usepackage[table,xcdraw,dvipsnames]{xcolor}
\usepackage{colortbl}
\usepackage{multirow}
\usepackage{makecell}
\usepackage{arydshln}

% 박스/밑줄/꾸미기
\tcbuselibrary{breakable}
\usepackage[normalem]{ulem}
\useunder{\uline}{\ul}{}
\usepackage{enumitem}
\usepackage{pifont}


% 알고리즘 (algorithmicx/algpseudocode 권장)
\usepackage{algorithm}
\usepackage{algorithmic}
% \usepackage{algpseudocode}
\newcommand{\ECT}{\textsc{ExtractConditionalTriplets}}
\newcommand{\forcomment}[1]{%
  \STATE \hfill
  {\color{green!60!black}\footnotesize$\triangleright$~#1}%
}

% 목차
\usepackage{hyperref}
\usepackage{titletoc}

% 프롬프트 템플릿 한페이지에 넣기 위함
\usepackage{cuted}
\usepackage{placeins}

% 수식
% Preamble 부분
\usepackage{amsmath}
% ------------------
% If the title and author information does not fit in the area allocated, uncomment the following
%
%\setlength\titlebox{<dim>}
%
% and set <dim> to something 5cm or larger.

% \title{ESGPT}
\title{ESGPT: A Data-Centric ESG Language Model with Unified Instruction and Benchmark Resources}

\author{
  \textbf{Jaeyoung Lee\textsuperscript{1}}
  \quad
  \textbf{Misuk Kim\textsuperscript{1,2,*}}
\\
  \textsuperscript{1}Department of Data Science, Hanyang University, Seoul, Republic of Korea
\\
  \textsuperscript{2}Department of Artificial Intelligence, Hanyang University, Seoul, Republic of Korea
\\
  \texttt{\{jaeylee, misukkim\}@hanyang.ac.kr}
\\
\textbf{GitHub:} \url{https://github.com/nowzer0/ESGPT}
}

\begin{document}
\maketitle

\renewcommand{\thefootnote}{}
\footnote{$^*$Corresponding author.}

\begin{abstract}
ESG analysis requires complex language understanding across diverse sources, including corporate reports, news, and disclosures, for tasks such as ESG relevance detection, action interpretation, impact analysis, and question answering. However, existing ESG datasets are fragmented by task and source, with heterogeneous formats and annotation schemes, limiting their use for training general-purpose ESG LLMs and unified evaluation. In this work, we construct \textsf{ESG-Instruct}, a 30,000-instance ESG instruction dataset built from public ESG resources, and \textsf{ESGU-Bench}, an ESG understanding benchmark covering six representative ESG tasks. Based on these resources, we further develop \textsf{ESGPT}, an 8B ESG-specialized LLM, and evaluate it against diverse open-source and API-based LLMs. Experimental results show that general-purpose LLMs exhibit substantial performance variation across ESG tasks, while \textsf{ESGPT}, despite its relatively small scale, achieves the best performance on multiple ESG tasks. This work provides unified data and benchmark resources for training and evaluating ESG-specialized LLMs.

% ESG 분석은 기업 보고서, 뉴스, 공시 등 다양한 출처의 정보를 기반으로 ESG 관련성 판단, 주제 식별, 행동 해석, 영향 분석, 질의응답 등 복합적인 언어 이해를 요구한다. 그러나 기존 ESG 데이터셋은 과제와 출처별로 분산되어 있으며 데이터 형식과 주석 체계가 상이하여, 범용적인 ESG LLM 학습과 통합 평가에 활용되기 어렵다. 본 연구는 공개 ESG 자원을 기반으로 30,000개 규모의 ESG instruction dataset인 \textsf{ESG-Instruct}와, 6개 대표 ESG task로 구성된 ESG understanding benchmark \textsf{ESGU-Bench}를 구축한다. 또한 이를 기반으로 8B 규모의 ESG 특화 LLM인 \textsf{ESGPT}를 개발하고, 다양한 오픈소스 및 API 기반 LLM과 비교 평가한다. 실험 결과, 범용 LLM들은 ESG 과제 간 큰 성능 편차를 보였으며, \textsf{ESGPT}는 상대적으로 작은 규모에도 불구하고 다수의 ESG 과제에서 최고 성능을 기록하였다. 본 연구는 ESG 특화 LLM의 학습과 평가를 위한 통합 데이터 및 벤치마크 자원을 제공한다. The dataset, benchmark, and model will be publicly released. To preserve anonymity during review, sample instances are provided via an anonymized repository: \url{ANONYMOUS_URL}.
\end{abstract}

\section{Introduction}
\begin{figure*}[t!]
    \centering
    \includegraphics[width=1\linewidth]{Figure/Fig1_Overall.pdf}
    \caption{Overall accuracy by model size on the \textsf{ESGU-Bench} across 35 LLMs.}
    \label{fig:overall_acc}
\end{figure*}
Environmental, Social, and Governance (ESG) has become a key criterion for assessing corporate sustainability and social responsibility. ESG information appears in diverse forms, including sustainability reports, disclosures, news articles, and international standards, and requires complex understanding and reasoning beyond simple classification, such as relevance detection, claim interpretation, impact analysis, and question answering. Large Language Models (LLMs) have shown the potential to address these ESG tasks through natural language interfaces~\cite{sun2026esg}. However, as ESG is a specialized domain whose interpretation depends on industry context, disclosure standards, and supporting evidence, strong performance on general benchmarks does not necessarily indicate ESG understanding. Therefore, ESG-specific data and benchmarks are needed to evaluate ESG understanding capabilities~\cite{he2025esgenius}.
% Environmental, Social, and Governance (ESG)는 기업의 지속가능성과 사회적 책임을 평가하기 위한 핵심 기준으로 활용되고 있다. ESG 정보는 지속가능성 보고서, 공시, 뉴스, 국제 기준 문서 등 다양한 형태로 존재하며, 단순 분류를 넘어 관련성 판단, 주장 해석, 영향 분석, 질의응답 등 복합적인 이해와 추론을 요구한다. Large Language Models (LLMs)는 이러한 ESG 과제를 자연어 기반으로 통합 처리할 수 있는 가능성을 보였으나~\cite{sun2026esg}, ESG는 산업별 맥락, 공시 기준, 근거 출처에 따라 해석이 달라지는 특수 도메인이므로 일반 벤치마크에서의 높은 성능이 ESG 이해 능력을 보장하지는 않는다. 따라서 ESG 도메인 특성을 반영한 데이터와 ESG-specific capability를 평가할 수 있는 벤치마크가 필요하다~\cite{he2025esgenius}.

Despite this need, publicly available ESG resources remain limited~\cite{ong2025towards, lee2023esginfo, lee2024esgkor}. Existing datasets are fragmented across languages, input structures, output formats, annotation units, and evaluation criteria, limiting their use for training general-purpose ESG LLMs and for unified evaluation. Moreover, ESG data construction often requires access to restricted source materials and substantial domain expertise. This highlights the need to systematically leverage high-quality public ESG resources and build an integrated data foundation for consistent training and evaluation across diverse ESG tasks.
% 이러한 필요성에도 불구하고 공개 ESG 데이터 자원은 여전히 제한적이며~\cite{ong2025towards, lee2023esginfo, lee2024esgkor}, 기존 데이터셋들은 언어, 입력 구조, 출력 형식, 주석 단위, 평가 기준이 서로 상이하여 범용 ESG LLM 학습과 통합 평가에 활용되기 어렵다. 또한 ESG 데이터 구축은 원천 소스 접근성 제약과 높은 전문가 의존도를 수반하므로, 공개된 고품질 ESG 자원을 체계적으로 활용하는 것이 중요하다. 따라서 다양한 ESG 과제를 일관된 방식으로 학습·평가할 수 있는 통합 데이터 기반이 요구된다.
This work constructs \textsf{ESG-Instruct}, a large-scale ESG instruction dataset, by systematically selecting and processing public ESG data resources. Based on this dataset, we develop \textsf{ESGPT}, a data-centric ESG LLM, and propose \textsf{ESGU-Bench}~(ESG Understanding Benchmark), which is based on human annotation and expert review. Experiments on 35 LLMs ranging from 0.5B to 2.2T parameters show that general-purpose LLMs exhibit substantial variation across ESG tasks, and even some large API-based LLMs show limited performance. In contrast, \textsf{ESGPT} achieves strong performance on multiple ESG tasks despite its relatively small scale. The contributions of this work are as follows:

\begin{itemize}
    \item We construct \textsf{ESG-Instruct}, a large-scale ESG instruction dataset, by systematically selecting and integrating public ESG data resources, enabling diverse ESG understanding tasks to be learned in a consistent instruction format.
    \item We construct \textsf{ESGU-Bench}, consisting of six representative tasks centered on human-annotated and expert-reviewed datasets, to reliably evaluate the ESG understanding capability of LLMs.
    \item We develop \textsf{ESGPT}, a data-centric ESG LLM based on \textsf{ESG-Instruct}, and analyze the limitations of general-purpose LLMs and the effectiveness of ESG-specialized models through large-scale comparative experiments.
\end{itemize}
% 본 연구는 공개 ESG 데이터 자원을 체계적으로 선별·가공하여 대규모 ESG instruction dataset인 \textsf{ESG-Instruct}를 구축한다. 이를 기반으로 data-centric ESG LLM인 \textsf{ESGPT}를 개발하며, 인간 주석과 전문가 검토 기반의 \textsf{ESGU-Bench}~(ESG Understanding Benchmark)를 제안한다. 0.5B--2.2T 규모의 35개 LLM에 대한 실험 결과, 범용 LLM은 ESG understanding에서 과제별 편차를 보였으며 일부 대형 API 기반 LLM 모델도 제한적인 성능을 나타냈다. 반면 \textsf{ESGPT}는 상대적으로 작은 규모에도 불구하고 다수의 ESG 과제에서 우수한 성능을 달성하였다. 본 연구의 기여는 다음과 같다:

% \begin{itemize}
%     \item 공개 ESG 데이터 자원을 체계적으로 선별·통합하여, 다양한 ESG understanding task를 일관된 instruction format으로 학습할 수 있는 대규모 ESG instruction dataset \textsf{ESG-Instruct}를 구축한다.
%     \item 인간 주석과 전문가 검토 기반 데이터셋을 중심으로, 6개 대표 과제로 구성된 \textsf{ESGU-Bench}를 구축하여 LLM의 ESG understanding capability를 신뢰성 있게 평가한다.
%     \item \textsf{ESG-Instruct}를 기반으로 data-centric ESG LLM인 \textsc{ESGPT}를 개발하고, 대규모 비교 실험을 통해 범용 LLM의 한계와 ESG 특화 모델의 효과를 분석한다.
% \end{itemize}

\section{Related Work}
Datasets and benchmarks for ESG analysis have been built from diverse sources, including corporate reports, news articles, disclosures, and climate documents. Early studies focused mainly on classification tasks such as ESG relevance detection, topic classification, and sentiment analysis, while recent work has expanded to more complex ESG understanding tasks, including corporate action interpretation, claim verification, and question answering. ESG-FTSE~\cite{pavlova2024esgftse} and ESG-Kor~\cite{lee2024esgkor} address ESG sentence detection and classification, whereas A3CG~\cite{ong2025towards} and PromiseEval~\cite{seki2025mlpromise} target fine-grained ESG interpretation of claims, actions, promises, and evidence. SustainableQA~\cite{ali2025sustainableqa} and ESGenius~\cite{he2025esgenius} further provide QA-based benchmarks for evaluating ESG and sustainability knowledge.
% ESG 분석을 위한 데이터셋과 벤치마크는 기업 보고서, 뉴스, 공시, 기후 문서 등 다양한 출처를 기반으로 구축되어 왔다. 초기 연구는 ESG 관련성 판단, 주제 분류, 감성 분석과 같은 분류 중심 과제에 초점을 맞추었으며, 최근에는 기업의 행동 해석, 주장 검증, 질의응답 등 보다 복합적인 ESG understanding task로 확장되고 있다. ESG-FTSE~\cite{pavlova2024esgftse}와 ESG-Kor~\cite{lee2024esgkor}는 ESG 문장의 탐지 및 분류를 다루며, A3CG~\cite{ong2025towards}와 PromiseEval~\cite{seki2025mlpromise}은 claim, action, promise, evidence와 같은 세밀한 ESG 해석 문제를 다룬다. 또한 SustainableQA~\cite{ali2025sustainableqa}와 ESGenius~\cite{he2025esgenius}는 ESG 및 sustainability knowledge 평가를 위한 QA 기반 benchmark로 활용될 수 있다.

From the model perspective, domain-specific pre-trained language models (PLMs), such as ESGBERT~\cite{schimanski2024ESGBERT} and ClimateBERT~\cite{webersinke2021climatebert}, have been used for ESG and climate-related text analysis. More recently, research has expanded toward LLM- or retrieval-augmented generation-based ESG analysis, question answering, and report generation, as in ChatReport~\cite{ni2023chatreport}, ESGReveal~\cite{zou2025esgreveal}, and SusGen-GPT~\cite{wu2025susgen}.
% 모델 측면에서는 ESGBERT~\cite{schimanski2024ESGBERT}, ClimateBERT~\cite{webersinke2021climatebert}와 같은 domain-specific PLM이 ESG 및 climate-related text 분석에 활용되었으며, 최근에는 ChatReport~\cite{ni2023chatreport}, ESGReveal~\cite{zou2025esgreveal}, SusGen-GPT~\cite{wu2025susgen}와 같이 LLM 또는 retrieval augmented generation 기반의 ESG 분석, 질의응답, 보고서 생성 등의 연구로 확장되고 있다.

However, existing studies mainly focus on individual tasks, specific sources, or particular applications. As a result, integrated instruction datasets for general-purpose ESG LLM development and reliable benchmarks for ESG understanding remain limited. Therefore, this work jointly presents an ESG instruction dataset, an ESG understanding benchmark, and an ESG-specialized LLM by consolidating and reformulating public ESG resources.
% 그러나 기존 연구들은 주로 개별 과제, 특정 출처, 또는 특정 응용에 초점을 맞추고 있어, 범용 ESG LLM 개발을 위한 통합 instruction dataset과 신뢰성 있는 ESG understanding benchmark는 여전히 부족하다. 따라서 본 연구는 공개 ESG 자원을 통합적으로 재구성하여 ESG instruction dataset, ESG understanding benchmark, 그리고 ESG 특화 LLM을 함께 제시한다.

\section{ESG Resource \& Model Construction}
\label{sec:construction}
This study collects and reorganizes public ESG datasets to build integrated resources for training and evaluating ESG-specialized LLMs. The data are organized by task objective, language, source type, annotation scheme, answer structure, and output requirements, followed by quality control to identify duplicates, missing values, formatting errors, label inconsistencies, and task-unsuitable instances. While preserving the original task definitions and answer structures, we convert the data into chat-prompt-based instances through translation, format normalization, task refinement, and prompt design. The refined tasks are then reorganized from the perspective of ESG understanding capabilities and integrated into six high-level categories. Dataset construction details are provided in Appendix~\ref{app:esg_resources}.

% 본 연구는 공개 ESG 데이터셋을 수집·재구성하여 ESG 특화 LLM 학습과 평가를 위한 통합 자원을 구축한다. 수집된 데이터는 과제 목적, 언어, 원천 데이터 유형, 주석 방식, 정답 체계, 출력 조건을 기준으로 체계화되며, 중복, 결측값, 형식 오류, 라벨 일관성, task 적합성 등을 직접 검토하는 품질 관리 과정을 거친다. 이후 원래의 task 정의와 정답 체계를 유지한 채 번역, 포맷 정규화, task 세분화, prompt 설계를 통해 chat-prompt 기반 인스턴스로 변환한다. 또한 세분화된 task를 ESG understanding capability 관점에서 재구성하여 6개 상위 카테고리로 통합한다. Dataset construction details are provided in Appendix~\ref{app:esg_resources}.
\begin{figure}[t]
    \centering
    \includegraphics[width=1\linewidth]{Figure/Fig2_main.pdf}
    \caption{Main results on \textsf{ESGU-Bench}: category-level accuracy (top) and overall accuracy of LLMs (bottom). Dashed vertical lines indicate the average overall accuracy of each model group.}
    \label{fig:main_results}
\end{figure}
\paragraph{\textsf{ESG-Instruct}.}
\textsf{ESG-Instruct} is a 30,000-instance instruction dataset for training ESG-specialized LLMs. From quality-reviewed ESG data, we construct 35 instruction tasks and select training instances by considering task coverage and data balance~\cite{wu2025susgen}. We adjust data proportions across the six ESG understanding categories to mitigate bias toward specific tasks or datasets. To support both ESG domain adaptation and general instruction-following ability, we allocate 10\% of the dataset to general instruction data. Considering the limited availability of public ESG resources, we also include translated data and a small amount of LLM-assisted data.

% \textsf{ESG-Instruct}는 ESG 특화 LLM 학습을 위한 30,000개 규모의 instruction dataset이다. 품질 검토를 통과한 ESG 데이터를 기반으로 총 35개의 instruction task를 구성하며, task coverage와 데이터 균형을 고려하여 학습 데이터를 선별한다~\cite{wu2025susgen}. 또한 6개 ESG 이해 카테고리를 기준으로 데이터 비율을 조정하여 특정 과제나 데이터셋의 편향을 완화한다. 마지막으로 ESG domain adaptation과 일반적인 instruction-following 능력을 함께 유지하기 위해 전체 데이터의 10\%를 general instruction data로 구성한다. 공개 ESG 자원의 제한성을 고려하여 번역 데이터와 일부 LLM-assisted data도 포함한다.

\paragraph{\textsf{ESGU-Bench}.}
\textsf{ESGU-Bench} is an ESG understanding benchmark composed of human-annotated and expert-reviewed datasets. It includes 9,514 instances and 21 evaluation tasks, which are integrated into six high-level ESG understanding categories to evaluate diverse ESG understanding capabilities beyond single classification- or QA-centered evaluation.
% \textsf{ESGU-Bench}는 인간 주석 및 전문가 검토 기반 데이터셋으로 구성된 ESG understanding benchmark이다. 총 9,514개의 instances와 21개 evaluation task를 포함하며, 이를 6개 상위 ESG 이해 카테고리로 통합하여 단일 분류 또는 QA 중심 평가를 넘어 다양한 ESG understanding capability를 평가한다.

\paragraph{\textsf{ESGPT}.}
\textsf{ESGPT} is a data-centric ESG LLM trained on \textsf{ESG-Instruct}. Rather than introducing a new model architecture, this study focuses on analyzing how the design of ESG instruction data, considering task coverage, source diversity, label balance, and data quality, affects ESG understanding performance.
% \textsf{ESGPT}는 \textsf{ESG-Instruct} 기반으로 학습한 data-centric ESG LLM이다. 본 연구는 새로운 모델 구조보다, task coverage, source diversity, label balance, data quality를 고려한 ESG instruction data 설계가 ESG understanding 성능에 미치는 영향을 분석하는 데 초점을 둔다.

\begin{figure}[t]
    \centering
    \includegraphics[width=1\linewidth]{Figure/Fig3_Ablation.pdf}
    \caption{Effect of \textsf{ESG-Instruct} data scale on \textsf{ESGPT}. Lines denote overall accuracy on \textsf{ESGU-Bench}, and bars denote training time in hours. All data-scale variants were trained with the same fixed random seed and training configuration and evaluated under the same protocol.}
    \label{fig:ablation}
\end{figure}

\section{Experiments}
\label{sec:experimental_setup}
Using \textsf{ESGU-Bench}, we conduct large-scale comparative experiments on 35 LLMs of diverse types and scales. All models are evaluated in a zero-shot setting, and performance is reported in terms of accuracy for each of the six ESG understanding categories and overall accuracy. All experiments were conducted on a single NVIDIA A100 GPU. Training and experimental details are provided in Appendix~\ref{app:implementation}, and detailed results and analyses are provided in Appendix~\ref{app:additional_results}.
% \textsf{ESGU-Bench}를 통해, 다양한 유형과 규모의 35개 LLMs을 대상으로 대규모 비교 실험을 수행한다. 모든 모델은 zero-shot setting에서 평가하며, 성능은 6개 ESG 이해 카테고리별 accuracy와 overall accuracy로 보고한다. 모든 실험은 단일 NVIDIA A100 GPU에서 수행되었다. 학습 및 실험 디테일은 Appendix~\ref{app:implementation}에서 제시하며, 세부 결과 및 분석 is in Appendix~\ref{app:additional_results}.

\paragraph{\textsf{ESGU-Bench} Results.}
Figure~\ref{fig:main_results} shows the performance of representative LLMs on \textsf{ESGU-Bench}. In the main text, we report results focusing on the model with the highest overall accuracy within each model family. ESG understanding performance varies substantially across model families and task types, and even large proprietary LLMs do not show consistent advantages across all ESG categories.
In contrast, \textsf{ESGPT}, despite being an 8B model, achieves the best performance in four of the six categories and outperforms both large API-based models from the DeepSeek and GPT families and open-source ESG LLMs, with an overall accuracy of 66.96\%. This demonstrates that systematically constructed ESG instruction data alone can lead to strong ESG understanding performance. At the same time, performance limitations are still observed in some ESG tasks, suggesting the need for more refined ESG contextual understanding and domain-specific reasoning. 

We also qualitatively analyze the ESG understanding performance of \textsf{ESGPT} and general-purpose LLMs. Input prompts and inputs are presented in summarized form.
% Figure~\ref{fig:main_results}는 \textsf{ESGU-Bench}에서 주요 LLM들의 성능을 보여준다. 본문에서는 각 모델 패밀리에서 가장 높은 overall accuracy를 보인 모델을 중심으로 결과를 제시한다. ESG understanding 성능은 모델 계열과 과제 유형에 따라 큰 차이를 보였으며, 대형 proprietary LLM도 모든 ESG 카테고리에서 일관된 우위를 보이지는 않았다. 
% 반면 \textsf{ESGPT}는 8B 규모의 모델임에도 6개 카테고리 중 4개에서 최고 성능을 기록하였으며, overall accuracy 66.96\%로 DeepSeek, GPT 계열의 대형 API 모델과 open-source ESG LLM을 모두 능가하였다. 이는 체계적으로 구성된 ESG instruction data만으로도 강한 ESG understanding 성능을 달성할 수 있음을 보여준다. 동시에 일부 ESG task에서는 여전히 성능 한계가 관찰되어, 보다 정교한 ESG 문맥 이해와 도메인 추론 능력이 필요함을 시사한다. 또한 \textsf{ESGPT}와 범용 LLM의 ESG 이해 성능을 정성적으로 분석한다. 입력 프롬프트 및 입력은 요약하여 제시한다. 
\begin{maincasebox}
\small
\textbf{[User]} 
Classify the ESG impact severity. \newline Allowed labels: Low, Medium, High.
\vspace{0.45em}

\textbf{[Input]} A.P. Møller acquires Faerch Group, a sustainable food packaging company providing circular food packaging solutions, products made from up to 100\% recycled content, fully recyclable food-grade packaging, and an integrated recycling system.
\tcblower
\scriptsize
\textbf{[Gold and Model Answers]}
\vspace{0.15em}

\setlength{\tabcolsep}{2.5pt}
\renewcommand{\arraystretch}{1.02}
\begin{tabularx}{1.0\linewidth}{@{}Y C{0.30\linewidth} C{0.30\linewidth}@{}}
\hline
\textbf{Model} & \textbf{Answer} & \textbf{Result} \\
\hline
Gold & \goldans{High} & -- \\
\hline
DeepSeek & \wrongans{Medium} & \xmarkred \\
GPT & \wrongans{Low} & \xmarkred \\
\textsf{ESGPT} & \correctans{High} & \cmarkblue \\
\hline
\end{tabularx}
\end{maincasebox}
\textsf{ESGPT} predicts \texttt{High} by recognizing the acquisition of Faerch Group, with its \textit{circular packaging and integrated recycling system}, as a \textit{strategically significant ESG event}. In contrast, the DeepSeek and GPT-family models appear to treat the case as \textit{a general ESG-related corporate activity} and assign a lower impact severity. This suggests that an ESG-specialized LLM can better capture the \textit{ESG significance and impact severity of corporate events}. 
Additional cases and detailed analyses are provided in Appendix~\ref{app:case_studies}.

% \textsf{ESGPT}는 \textit{순환형 포장 솔루션, 재활용 소재, 완전 재활용 가능성, 통합 재활용 시스템 등 높은 ESG impact를 시사하는 근거}를 반영하여 \texttt{High}를 예측하였다. 반면 DeepSeek와 GPT 계열 모델은 이를 \textit{일반적인 ESG 관련 기업 활동으로 해석하여 impact severity를 낮게 평가}하였다. 이는 ESG 특화 LLM이 단순한 ESG 관련성 뿐만 아니라, 사건의 ESG적 중요도와 영향 수준을 더 잘 포착할 수 있음을 보여준다. 
% 추가 사례와 상세 분석은 Appendix~\ref{app:case_studies}에 제시한다. 

\paragraph{\textsf{ESG-Instruct} Scale Analysis.}
Figure~\ref{fig:ablation} shows performance changes according to the scale of \textsf{ESG-Instruct}. With only 10K training instances, the model already ranks among the top-performing models, and further improves when scaled to 30K. This trend is consistent across different backbone models. By contrast, the full-scale setting degrades performance despite requiring the highest training cost. These results suggest that task/category balance and data quality are more important for ESG instruction tuning than simply increasing data scale. In particular, tasks with limited high-quality annotations are difficult to retain at sufficient proportions in the full-scale setting, while data-rich tasks may become overrepresented.
% Figure~\ref{fig:ablation}은 \textsf{ESG-Instruct} 규모에 따른 성능 변화를 보여준다. 10K 규모의 데이터만으로도 전체 LLM 중 second-best 수준의 성능을 달성하였으며, 30K까지 확장할 때 성능이 추가적으로 향상되었다. 이러한 경향은 서로 다른 backbone model에서도 일관되게 관찰되었다. 반면 full-scale setting에서는 가장 높은 학습 비용에도 불구하고 성능이 감소하였다. 이는 단순한 데이터 규모 확장보다, task/category 균형과 데이터 품질을 고려한 구성 전략이 ESG instruction tuning에서 더 중요함을 보여준다. 특히 고품질 주석 데이터가 제한적인 과제는 full-scale 세팅에서 충분한 비율을 유지하기 어렵고, 데이터 규모가 큰 특정 task가 과도하게 반영될 수 있다.

\begin{table}[t]
\centering
\small
\setlength{\tabcolsep}{3pt}
\renewcommand{\arraystretch}{1.12}
\begin{tabularx}{\linewidth}{@{}XX@{}}
\hline
\textbf{\textsf{ESGPT}} & \textbf{SusGen-GPT} \\
\hline
The board indirectly oversees climate-related risks through operational safety and financial leverage reduction, but \textcolor{caseblue}{the report does not explicitly link these strategies to climate-specific measures}. 
&
The board oversees climate-related risks through governance mechanisms, and the Petrobras Corruption Prevention Program \textcolor{casered}{likely includes climate-related risks}. \\
\hline
\end{tabularx}
\caption{Compressed example of TCFD-style report generation.}
\label{tab:main_report_generation_case}
\end{table}

\paragraph{Report Generation Analysis.}
We qualitatively compare the TCFD~(Task Force on Climate-related Financial Disclosures) report generation behavior of \textsf{ESGPT} and SusGen-GPT, a sustainability report generation model (Table~\ref{tab:main_report_generation_case}). A representative case is presented in the main text, with the full comparison in Appendix~\ref{app:report_generation}.

\textsf{ESGPT} notes that operational safety and financial strategy may be indirectly related to climate risks, but does not infer climate-specific oversight without explicit evidence. In contrast, SusGen-GPT links even the Petrobras Corruption Prevention Program to climate-related risk oversight, extending general governance information into climate-specific disclosure. Since ESG report generation is tied to corporate disclosure, such unsupported extensions can undermine reliability. This case suggests that \textsf{ESGPT} better preserves the boundary between general ESG information and climate-specific disclosure, supporting more reliable and evidence-grounded ESG report generation.

% \textsf{ESGPT}와 지속가능성 보고서 생성 특화 모델인 SusGen-GPT의 TCFD~(Task Force on Climate-related
% Financial Disclosures) report 생성 양상을 정성적으로 비교한다 (Table~\ref{tab:main_report_generation_case}). 본문에서는 대표 사례로 제시하며, 전체 비교는 Appendix~\ref{app:report_generation}에 제시한다.

% \textsf{ESGPT}는 운영 안전성 및 재무 전략이 기후 리스크와 간접적으로 관련될 수 있음을 언급하면서도, climate-specific oversight가 명시되지 않았다는 한계를 유지한다. 반면 SusGen-GPT는 Petrobras Corruption Prevention Program까지 climate-related risk oversight로 연결하며 일반 governance 정보를 더 강하게 확장한다. ESG 보고서 생성은 기업 공시와 연결되는 작업이므로, 근거가 명확하지 않은 내용을 기후 관련 공시로 확장하는 것은 보고서 신뢰성을 저해할 수 있다. 이 사례는 \textsf{ESGPT}가 일반 ESG 정보와 climate-specific disclosure의 경계를 더 안정적으로 구분하며, 신뢰도 높은 보고서 생성에 강점을 가질 수 있음을 보여준다.

\section{Conclusion}
\label{sec:conclusion}
This study proposed \textsf{ESG-Instruct}, \textsf{ESGU-Bench}, and \textsf{ESGPT}, a data-centric ESG LLM, for ESG-specialized LLM development and reliable evaluation. These resources are designed to support consistent training and evaluation across diverse ESG understanding tasks. Experiments show that general-purpose LLMs and large API-based models exhibit substantial performance variation across ESG tasks, whereas \textsf{ESGPT} achieves strong ESG understanding performance despite its relatively small scale. These results demonstrate the effectiveness of the proposed ESG domain-specific instruction data and integrated benchmark for training and evaluating ESG-specialized LLMs.
% 본 연구는 ESG 특화 LLM 개발과 신뢰성 있는 평가를 위해 ESG instruction dataset \textsf{ESG-Instruct}, ESG understanding benchmark \textsf{ESGU-Bench}, 그리고 data-centric ESG LLM인 \textsf{ESGPT}를 제안하였다. 제안한 자원은 다양한 ESG understanding task를 일관된 방식으로 학습하고 평가할 수 있도록 설계되었다. 실험 결과, 범용 LLM과 대형 API 기반 모델은 ESG 과제 간 큰 성능 편차를 보였으며, \textsf{ESGPT}는 상대적으로 작은 규모에도 불구하고 우수한 ESG understanding 성능을 달성하였다. 이는 본 논문에서 제안한 ESG domain-specific instruction data와 통합 benchmark가 ESG 특화 LLM 학습 및 평가에 효과적임을 보여준다.

\section*{Limitations}
Although this study proposes an integrated benchmark for ESG understanding and an ESG-specialized instruction model, it has several limitations. 
First, while \textsf{ESGU-Bench} covers diverse ESG understanding tasks, including ESG relevance detection, topic identification, commitment analysis, impact analysis, claim verification, and question answering, it cannot fully represent all ESG disclosure contexts across industries, regions, and regulatory frameworks. 
ESG issues can be interpreted differently depending on industry-specific materiality and country-specific disclosure requirements. 
Nevertheless, we collected a wide range of publicly available ESG datasets and constructed the benchmark while considering the distribution of tasks and data sources, in order to maintain broad coverage of ESG understanding as much as possible.
% 본 연구는 ESG understanding을 위한 통합 benchmark와 ESG 특화 instruction model을 제안하였으나, 몇 가지 한계를 가진다. 
% 첫째, \textsf{ESGU-Bench}는 ESG relevance detection, topic identification, commitment analysis, impact analysis, claim verification, reasoning 등 다양한 ESG understanding task를 포함하지만, 모든 산업, 지역, 규제 프레임워크의 ESG 공시 환경을 완전히 포괄하지는 못한다. ESG 이슈는 산업별 중요도와 국가별 공시 기준에 따라 다르게 해석될 수 있기 때문이다. 
% 그럼에도 본 연구는 공개적으로 활용 가능한 다양한 ESG 데이터셋을 수집하고, 특정 task나 데이터 출처가 과도하게 반영되지 않도록 구성함으로써 가능한 한 폭넓은 ESG understanding 범위를 유지하고자 하였다.

Second, the ESG data used in this study is largely English-centric. 
This reflects the fact that many publicly available ESG reports, news articles, standard documents, and benchmark datasets are primarily constructed in English. 
We translated some non-English ESG data into English and used them for training; however, regional expressions, institutional nuances, and disclosure practices in the original texts may have been partially weakened during translation. 
Nevertheless, during the translation process, we sought to preserve the original meaning as much as possible by retaining numerical values, proper nouns, commitment expressions, verification-related wording, paragraph structure, and the domain-specific context of ESG, sustainability, finance, and disclosure. Future work should extend \textsf{ESG-Instruct} and \textsf{ESGU-Bench} toward multilingual and jurisdiction-aware ESG resources, and further develop \textsf{ESGPT} to better reflect such linguistic and regional diversity. Since high-quality ESG resources remain limited for non-English languages and region-specific disclosure contexts, future ESG-specialized datasets, benchmarks, and models should better capture differences in regulatory environments, regional reporting practices, institutional terminology, and local ESG expressions.
% 둘째, 본 연구에서 사용한 ESG 데이터는 영어 자료의 비중이 높다. 
% 이는 공개적으로 활용 가능한 ESG 보고서, 뉴스, 표준 문서, benchmark 데이터가 주로 영어 중심으로 구축되어 있기 때문이다. 
% 본 연구에서는 일부 비영어권 ESG 데이터를 영어로 번역하여 학습 데이터로 활용하였으나, 번역 과정에서 원문의 지역적 표현, 제도적 뉘앙스, 공시 관행의 차이가 일부 약화될 가능성이 있다. 
% 그럼에도 번역 과정에서 ESG, 지속가능성, 금융, 공시 문맥의 의미가 훼손되지 않도록 원문의 수치, 고유명사, 약속 표현, 검증 표현, 문단 구조를 최대한 보존하고자 하였다.

Third, due to training resource constraints, \textsf{ESGPT} was developed by applying LoRA-based training to an 8B-scale base model. 
This setting enables ESG-specialized adaptation at a relatively limited training cost, but this study does not fully examine how performance may change with larger-scale models or full fine-tuning.
Our goal is not to exhaustively optimize architecture scaling, but to provide \textsf{ESG-Instruct}, \textsf{ESGU-Bench}, and \textsf{ESGPT} as a unified foundation for ESG-specialized LLM research and to demonstrate the feasibility of resource-efficient ESG adaptation.
Since ESG understanding requires fine-grained label boundary recognition, disclosure-context understanding, and evidence-based judgment, the impact of model scale and training strategy should be further investigated in future work.

Finally, ESG-specialized models may produce incorrect or incomplete judgments, especially in cases that require industry-specific expertise, regulatory interpretation, or evidence-based verification. Therefore, \textsf{ESGPT} and \textsf{ESGU-Bench} should be used as decision-support resources rather than substitutes for expert review in corporate disclosure, investment, or regulatory decision-making. 
Despite these limitations, this study provides a systematic benchmark for ESG understanding and demonstrates the potential of ESG-specialized LLMs, offering a foundation for developing more reliable ESG AI systems.
% 셋째, \textsf{ESGPT}는 학습 리소스 제약을 고려하여 8B 규모의 base model을 LoRA 기반으로 학습한 모델이다. 
% 이는 제한된 비용으로 ESG 특화 성능을 확보할 수 있다는 장점이 있지만, 더 큰 규모의 모델이나 full fine-tuning 환경에서의 성능 변화는 본 연구에서 충분히 검토하지 못하였다. 
% 특히 ESG understanding은 세밀한 label boundary, 공시 문맥 이해, 근거 기반 판단을 요구하므로, 모델 규모와 학습 방식에 따른 성능 변화는 향후 연구에서 추가적으로 분석될 필요가 있다. 
% 이러한 한계에도 불구하고, 본 연구는 ESG understanding을 체계적으로 평가하기 위한 benchmark와 ESG 특화 LLM의 가능성을 제시하며, 향후 신뢰 가능한 ESG AI 시스템 개발을 위한 기반을 제공한다.

\section*{Ethical Considerations}
This study uses publicly available ESG datasets and does not involve human subject experiments or personal information. The original dataset creators and sources are credited in Appendix~\ref{app:esg_resources}. Therefore, this work does not raise additional ethical concerns related to human subjects research or personal data processing.

\section*{Acknowledgments}
This work was supported by the Institute of Information \& Communications Technology Planning \& Evaluation (IITP) grant funded by the Korea government (MSIT) (No. RS-2020-II201373, Artificial Intelligence Graduate School Program (Hanyang University)) and the Institute of Information \& Communications Technology Planning \& Evaluation (IITP) grant funded by the Korea government (MSIT, Ministry of Science and ICT) (No. RS-2026-25548925).

% Custom bibliography entries only
\bibliography{reference}

\clearpage
\appendix
\setcounter{table}{0}
\setcounter{figure}{0}
\setcounter{equation}{0}

\renewcommand{\thetable}{A.\arabic{table}}
\renewcommand{\thefigure}{A.\arabic{figure}}

% ---------------------------------------------
\appendix

\section{ESG Resources}
\label{app:esg_resources}
We describe the collection and construction process of the ESG data resources used in this study. We collected ESG datasets with clear provenance and reproducibility for training and evaluating ESG-specialized LLMs. Datasets were selected primarily from public repositories such as Hugging Face, GitHub, and Kaggle, or from resources whose construction process, annotation scheme, label system, and source materials could be verified through papers or official documents. Some datasets were also obtained by directly requesting them from the original authors or providers. In such cases, we included them only when their provenance, task definition, annotation scheme, and answer format were verifiable. This process excludes arbitrarily collected or unofficial data whose provenance or answer structure cannot be verified.

The collected datasets cover diverse ESG-related sources, including corporate sustainability reports, climate disclosures, news and media articles, and ESG standards and frameworks. For each dataset, we manually reviewed its task objective, language, source type, annotation scheme, label system, answer format, and output requirements to assess its suitability as a training or evaluation resource. Table~\ref{tab:esg_resources} presents the datasets used in this study and their related information.

Since the collected datasets have different input structures and output formats, we applied a common preprocessing process for LLM training and evaluation. Specifically, we removed duplicate samples, checked missing fields, corrected formatting errors, normalized label names, unified answer formats, and reviewed task suitability. When a dataset contained multiple subtasks or hierarchical labels, we refined the task units so that each task objective and output requirement were clearly specified.

Through this process, we constructed two ESG data resources. First, \textsf{ESG-Instruct} is an instruction-tuning dataset for training ESG-specialized LLMs, aiming to support domain adaptation across diverse ESG understanding tasks. Second, \textsf{ESGU-Bench} is a benchmark for reliably evaluating the ESG understanding capability of LLMs, centered on evaluation data with human annotations, expert review, or clear gold answers.
% 본 연구에서 활용한 ESG 데이터 자원의 수집 및 구성 과정을 상세히 설명한다. 본 연구는 ESG 특화 LLM 학습과 평가에 활용하기 위해 출처가 명확하고 재현 가능한 ESG 데이터셋을 수집하였다. 데이터셋은 Hugging Face, GitHub, Kaggle 등 공개 저장소를 통해 배포되었거나, 논문 및 공식 문서를 통해 데이터 구축 과정, 주석 방식, 라벨 체계, 원천 자료를 확인할 수 있는 경우를 중심으로 선별하였다. 또한 일부 데이터셋은 원 저자 또는 제공 기관에 직접 요청하여 제공받았으며, 이 경우에도 데이터의 출처, 과제 정의, 주석 체계, 정답 형식이 확인 가능한 경우에 한해 본 연구의 데이터 자원으로 포함하였다. 이를 통해 출처와 정답 체계가 검증 가능하지 않은 임의 수집 데이터나 비공식 데이터는 제외하였다.

% 수집된 데이터셋은 기업 지속가능성 보고서, 기후 공시 자료, 뉴스 및 미디어 기사, ESG 표준 및 프레임워크 등 다양한 ESG 관련 출처를 포함한다. 
% 본 연구는 각 데이터셋의 과제 목적, 언어, 원천 데이터 유형, 주석 방식, 라벨 체계, 정답 형식, 출력 조건 등을 종합적으로 manually 검토하여 학습 및 평가 자원으로의 활용성을 판단하였다. Table~\ref{tab:esg_resources}에는 본 연구에 사용된 데이터셋 목록 및 관련 정보를 제시한다.

% 수집된 데이터셋은 서로 다른 입력 구조와 출력 형식을 가지므로, LLM 학습 및 평가에 활용하기 위해 공통적인 전처리 과정을 수행하였다. 
% 구체적으로 중복 샘플 제거, 결측 필드 확인, 형식 오류 수정, 라벨명 정규화, 정답 형식 통일, task 적합성 검토를 진행하였다. 
% 또한 하나의 데이터셋 안에 여러 세부 과제나 계층적 라벨이 포함된 경우, 각 과제의 목적과 출력 조건이 명확히 드러나도록 task 단위를 세분화하였다.

% 이러한 과정을 통해 본 연구는 두 가지 ESG 데이터 자원을 구축하였다. 첫째, \textsf{ESG-Instruct}는 ESG 특화 LLM 학습을 위한 instruction tuning 데이터셋으로, 다양한 ESG 이해 과제에 대한 도메인 적응을 목표로 한다. 둘째, \textsf{ESGU-Bench}는 LLM의 ESG understanding 능력을 신뢰성 있게 평가하기 위한 benchmark로, 인간 주석, 전문가 검토, 또는 명확한 gold answer를 포함하는 평가 데이터를 중심으로 구성한다.
\begin{figure}[h]
    \centering
    \includegraphics[width=1\linewidth]{Figure/ApFig1_Instruct_data.pdf}
    \caption{Category distribution of \textsf{ESG-Instruct} 30K.}
    \label{fig:Dist_ESG_instruct}
\end{figure}
% ---------------------------------------------
\subsection{\textsf{ESG-Instruct}}
\label{app:esg_instruct_details}
The construction of \textsf{ESG-Instruct} proceeds in three main stages: translation, instruction conversion, and final training-data selection. Since high-quality data publicly available in the ESG domain are limited, we include not only human-annotated data but also multilingual ESG data expanded through translation and a limited amount of data that involved LLM assistance during construction. However, these data are used only for training ESG domain adaptation and are kept separate from \textsf{ESGU-Bench}, which is designed for reliable model evaluation.
\begin{table}[t]
\centering
\resizebox{\linewidth}{!}{%
\begin{tabular}{llrr}
\toprule
\textbf{Dataset} & \textbf{Translation Direction} & \textbf{Samples} & \textbf{Time} \\
\midrule
ESG-Kor & Korean $\rightarrow$ English & 118,946 & 130:17:04 \\
Dynamic ESG & Chinese $\rightarrow$ English & 2,200 & 6:24:04 \\
ML-ESG1 & French $\rightarrow$ English & 1,500 & 4:27:33 \\
ML-ESG2 & French $\rightarrow$ English & 1,018 & 3:04:08 \\
ML-ESG2 & Japanese $\rightarrow$ English & 1,121 & 3:01:23 \\
ML-ESG3 & French $\rightarrow$ English & 807 & 2:24:45 \\
ML-ESG3 & Japanese $\rightarrow$ English & 1,553 & 4:14:13 \\
ML-ESG3 & Korean $\rightarrow$ English & 1,000 & 2:53:58 \\
PromiseEval & French $\rightarrow$ English & 800 & 1:54:26 \\
PromiseEval & Japanese $\rightarrow$ English & 800 & 1:56:57 \\
\midrule
\multicolumn{2}{l}{\textbf{Total}} & \textbf{129,745} & \textbf{160:38:31} \\
\bottomrule
\end{tabular}%
}
\caption{Translation time for constructing \textsf{ESG-Instruct} from multilingual ESG datasets.}
\label{tab:translation_time}
\end{table}
%------------------------------------
\begin{table*}[t]
\centering
\small
\resizebox{\textwidth}{!}{%
\begin{tabular}{p{0.19\linewidth}p{0.68\linewidth}r}
\toprule
Check & Description & Overlap \\
\midrule
Exact input & Exact match over normalized inputs. & 0 \\
Input--target pair & Exact match over normalized input--target pairs. & 0 \\
Near-duplicate (3-gram) & 3-gram Jaccard similarity $\geq 0.80$ within each ESG category. & 0 \\
Near-duplicate (5-gram) & 5-gram Jaccard similarity $\geq 0.80$ within each ESG category. & 0 \\
Source/document & Exact match over available source/document identifiers after filtering generic identifiers. & 0 \\
\bottomrule
\end{tabular}
}
\caption{Contamination checks between \textsf{ESG-Instruct} and \textsf{ESGU-Bench}.}
\label{tab:contamination_check}
\end{table*}
First, non-English ESG datasets are translated so that they can be used for English-based training. We use the \texttt{Gemini 3.1 Flash-Lite} for translation and design a dedicated translation prompt to preserve numbers, named entities, temporal expressions, commitment expressions, and verification-related wording that are important in ESG, sustainability, finance, and corporate reporting contexts. In particular, the prompt instructs the model not to summarize or expand the original meaning and to preserve ambiguity when it exists in the source text. In total, we translate 129,745 non-English ESG samples. The overall processing time is 160 hours, 38 minutes, and 31 seconds, with a total translation cost of approximately \$187. The data scale and processing time for translation are reported in Table~\ref{tab:translation_time}, and an example of the translation prompt is provided below.

% \textsf{ESG-Instruct} 구축 과정은 크게 번역, instruction 변환, 그리고 최종 학습 데이터 선별의 순서로 진행된다. 
% ESG 도메인은 공개적으로 활용 가능한 고품질 instruction 데이터가 제한적이기 때문에, 본 연구에서는 인간 주석 기반 데이터뿐만 아니라 번역을 통해 확장한 다국어 ESG 데이터와 데이터 구축 과정에서 LLM이 일부 개입된 데이터도 학습 자원으로 제한적으로 포함하였다. 
% 다만 이러한 데이터는 ESG 도메인 적응을 위한 학습 목적으로만 사용하였으며, 모델의 신뢰성 있는 평가를 위한 \textsf{ESGU-Bench}와는 분리하였다.

% 먼저 비영어 ESG 데이터셋은 영어 기반 instruction tuning에 활용할 수 있도록 번역하였다. 
% 번역에는 Gemini 3.1 Flash-Lite 모델을 사용하였으며, 번역 과정에서는 ESG, 지속가능성, 금융 및 기업 보고서 문맥에서 중요한 수치, 고유명사, 시점, 약속 표현, 검증 표현이 보존되도록 별도의 번역 프롬프트를 설계하였다. 특히 원문의 의미를 요약하거나 확장하지 않도록 지시하고, 모호성이 존재하는 경우 이를 번역 결과에서도 유지하도록 하였다. 총 129,745개의 비영어 ESG 샘플을 번역하였으며, 전체 처리 시간은 160시간 38분 31초, 총 번역 비용은 약 \$187이다. 
% 번역에 사용된 데이터 규모와 처리 시간은 Table~\ref{tab:translation_time}에 제시하며, 사용한 번역 프롬프트 예시는 아래에 제시한다.

%------------------------------------
\begin{center}
\begin{tcolorbox}[
    breakable,
    enhanced jigsaw,
    width=0.95\linewidth,
    colback=white,
    colframe=black,
    coltitle=white,
    colbacktitle=black,
    title=\textbf{Translation Prompt},
    fonttitle=\bfseries,
    boxrule=0.8pt,
    arc=2mm,
    left=6pt,
    right=6pt,
    top=6pt,
    bottom=6pt,
    before skip=6pt,
    after skip=6pt
]
\small
\textbf{Instruction:}\\
You are a professional translator for ESG, sustainability, finance, and business reports.

\vspace{0.5em}
Translate the following French ESG report text into fluent English.

\vspace{0.5em}
\textbf{Rules:}
\begin{itemize}
    \setlength{\itemsep}{0pt}
    \setlength{\leftmargin}{1.2em}
    \item Preserve meaning exactly. Do not summarize, omit, or add information.
    \item Preserve reporting nuance, commitments, verification wording, dates, figures, percentages, and proper nouns.
    \item Keep ambiguity if present.
    \item Keep paragraph breaks if they exist.
    \item Output only the translation.
    \item Do not add notes, comments, explanations, or quotation marks.
\end{itemize}

\textbf{Source Text:}\\
\texttt{\{text\}}

\vspace{0.5em}
\textbf{Output:}\\
English translation
\end{tcolorbox}
\end{center}
Next, each dataset is converted into a chat-prompt-based instruction format. Each instance consists of a system prompt, task-specific instruction, an input text or evidence context, and a target answer. The target answer preserves the original gold label or gold answer so that the label space and answer structure defined in each dataset remain unchanged. Prompt examples used for each task are provided in Appendix~\ref{app:task_taxonomy_prompt}.

Finally, the training data are selected by considering the distributions across the six high-level ESG understanding categories and individual tasks, so that no particular dataset or task is overrepresented. Within each task, we review the label distribution, the alignment between the input and target answer, and the presence of formatting errors. Samples with unclear labels or those that do not match the task objective are excluded. Through this process, we construct the final 30K-scale \textsf{ESG-Instruct}, consisting of 35 ESG instruction tasks and general instruction data~\cite{taori2023alpaca}. The task-level composition of \textsf{ESG-Instruct} and the final number of selected instances are presented in Table~\ref{tab:esg_instruct_composition}. Figure~\ref{fig:Dist_ESG_instruct} shows the category-level distribution of the final 30K training data, and Figure~\ref{fig:Token_Dist_ESG_instruct} shows the token-level distribution of the training instances.
% 이후 각 데이터셋을 chat-prompt 기반 instruction 형식으로 변환하였다. 
% 각 인스턴스는 task-specific instruction, 입력 텍스트 또는 근거 문맥, target answer로 구성된다. 
% 이때 target answer는 gold label 또는 gold answer를 유지하여, 데이터셋마다 정의된 label space와 정답 체계가 유지되도록 구성하였다. 
% 각 task에 사용한 prompt 예시는 Appendix~\ref{app:task_taxonomy_prompt}에 제시한다.

% 최종 학습 데이터는 특정 데이터셋이나 특정 task가 과도하게 반영되지 않도록 6개의 상위 ESG understanding category와 개별 task 단위의 분포를 고려하여 선별하였다. 각 task 내부에서는 label distribution, 입력과 정답 간 대응성, 형식 오류 여부를 검토하였으며, 라벨이 불명확하거나 task 목적에 부합하지 않는 샘플은 제외하였다. 이러한 과정을 통해 최종적으로 35개의 instruction task로 구성된 30K 규모의 \textsf{ESG-Instruct}를 구축하였다. 
% 전체 \textsf{ESG-Instruct}의 task별 구성과 최종 선별 instance 수는 Table~\ref{tab:esg_instruct_composition}에 제시한다. 
% 또한 Figure~\ref{fig:Dist_ESG_instruct}는 최종 30K 학습 데이터의 category-level 분포를, Figure~\ref{fig:Token_Dist_ESG_instruct}는 학습 인스턴스의 token-level 분포를 보여준다.
%------------------------------------
\begin{figure*}[t]
    \centering
    \includegraphics[width=1\linewidth]{Figure/ApFig2_Instruct_token_dist.pdf}
    \caption{Token length distributions of the \textsf{ESG-Instruct}.}
    \label{fig:Token_Dist_ESG_instruct}
\end{figure*}
%------------------------------------
\begin{figure*}[t]
    \centering
    \includegraphics[width=1\linewidth]{Figure/ApFig3_Bench_token_dist.pdf}
    \caption{Token length distributions of the \textsf{ESGU-Bench}.}
    \label{fig:Token_Dist_ESGU_bench}
\end{figure*}
% ---------------------------------------------
\subsection{\textsf{ESGU-Bench}}
\label{app:esgu_details}
\textsf{ESGU-Bench} is constructed with priority on the reliability and comparability of evaluation data. To this end, the benchmark primarily includes datasets with human annotations, expert review, or clearly provided gold answers. In contrast, datasets whose construction mainly involved LLMs or whose label verification process could not be sufficiently confirmed are excluded from benchmark evaluation.

To ensure reproducibility and fair comparison, we use the official test split provided by the original dataset as benchmark instances whenever available. When no official test split is provided, we either use the subset released for evaluation in the original dataset or construct an evaluation subset through data quality review. Most benchmark tasks are converted into a chat-prompt format for model evaluation and are designed in a controlled manner for each task to improve consistency across evaluation conditions. For ESG Question Answering, we retain the original data format to preserve the given question, context, and answer structure~\cite{he2025esgenius}. The task-level composition of \textsf{ESGU-Bench} and the final number of evaluation instances are presented in Table~\ref{tab:esg_benchmark_composition}. Figure~\ref{fig:Token_Dist_ESGU_bench} shows the category-wise token length distribution of benchmark instances.

% \textsf{ESGU-Bench}는 평가 데이터의 신뢰성과 비교 가능성을 우선하여 구성하였다. 
% 이를 위해 benchmark에는 인간 주석, 전문가 검토, 그리고 명확하게 제공된 gold answer를 포함하는 데이터셋을 중심으로 선별하였다. 
% 반대로 데이터 구축 과정에서 LLM이 주로 개입하였거나, 라벨에 대한 충분한 검증 절차가 확인되지 않는 데이터셋은 benchmark 평가에서 제외하였다.
% 평가 결과의 재현성과 공정한 비교 가능성을 확보하기 위해, 가능한 경우 원 데이터셋에서 제공하는 official test split을 benchmark instance로 사용하였다. 별도의 official test split이 제공되지 않는 경우에는 원 데이터셋에서 평가용으로 제공된 subset을 사용하거나, 데이터 품질 검토를 거쳐 평가용 subset을 구성하였다. 모든 benchmark task는 모델 평가를 위해 chat-prompt 형식으로 변환되며, 평가 조건의 일관성을 높이기 위해 task별로 통제된 형태로 설계하였다. 전체 \textsf{ESGU-Bench}의 task별 구성과 최종 evaluation instance 수는 Table~\ref{tab:esg_benchmark_composition}에 제시하며, Figure~\ref{fig:Token_Dist_ESGU_bench}는 benchmark instance의 category-wise token length distribution을 보여준다.
% ---------------------------------------------
\subsection{Data Leakage Control}
\label{app:data_leakage}
We carefully separated \textsf{ESG-Instruct} and \textsf{ESGU-Bench} during data construction to prevent train--evaluation contamination. 
Although some dataset families, such as PromiseEval and the ML-ESG series, appear in both resources, this corresponds to dataset-family overlap rather than sample-level leakage. 
In particular, for multilingual datasets, non-English subsets were translated and used for \textsf{ESG-Instruct}, whereas English gold/test subsets were used for \textsf{ESGU-Bench}. 
Thus, belonging to the same dataset family does not imply that the same sample pool is shared between training and evaluation.

Because publicly available high-quality ESG training data are limited, \textsf{ESG-Instruct} includes translated data and limited LLM-assisted data to cover diverse ESG tasks, languages, and source types. 
In contrast, \textsf{ESGU-Bench} was constructed separately for reliable evaluation, focusing on instances with human annotations, expert review, or clear gold answers. 
Whenever available, we used the official test split of the original dataset.

To further verify the independence between the two resources, we conducted contamination checks between the final \textsf{ESG-Instruct} training instances and \textsf{ESGU-Bench} evaluation instances~(Table~\ref{tab:contamination_check}). 
Since non-English training instances were translated into English before training, potential cross-lingual overlaps were also checked in the same English text space through exact and n-gram near-duplicate matching.

% ---------------------------------------------
\section{ESG Task Taxonomy and Prompt Templates}
\label{app:task_taxonomy_prompt}
We provide the ESG task taxonomy and prompt examples used in this study. While preserving the task definitions and answer structures of diverse ESG datasets, we organize related tasks into six high-level ESG task categories for training-data construction and benchmark evaluation. Since each task differs in input structure, label space, answer format, and output requirements, we design a task-specific prompt template for each task. However, presenting prompts for all tasks would substantially increase the paper length; therefore, we provide representative prompt examples for each of the six high-level ESG task categories.
% 본 연구에서 사용한 ESG task taxonomy와 prompt 예시를 제공한다. 본 연구는 다양한 ESG 데이터셋의 과제 정의와 정답 체계는 유지하되, 학습 데이터 구성과 benchmark 평가를 위해 관련 task를 6개의 상위 ESG task category로 정리하였다.
% 각 task는 원 데이터셋의 입력 구조, label space, 정답 형식, 출력 조건이 서로 다르기 때문에, 본 연구에서는 task별 특성을 고려하여 개별 prompt template을 설계하였다. 
% 다만 모든 task의 prompt를 제시할 경우 지면이 과도하게 증가하므로, 6개 상위 ESG task category별 대표 prompt 예시만을 제시한다.
% ---------------------------------------------
\subsection{ESG Relevance Detection}
ESG relevance detection is the task of determining whether a given sentence- or document-level text is related to ESG issues. This category includes ESG relevance classification, climate-related text detection, and environmental claim detection. The model must determine whether the input text is substantively connected to ESG-related topics such as environmental, social, governance, climate, sustainability, and corporate responsibility issues.
% ESG relevance detection은 주어진 문장 또는 문서 단위 텍스트가 ESG 이슈와 관련되는지를 판단하는 task이다. 
% 이 category에는 ESG relevance classification, climate-related text detection, environmental claim detection 등이 포함된다. 
% 모델은 입력 텍스트가 환경, 사회, 거버넌스, 기후, 지속가능성, 기업 책임 등 ESG 관련 주제와 실질적으로 연결되는지를 판단해야 한다.

\begin{promptbox}{ESG Relevance Classification}
\small
\textbf{[System]}\\
You are an ESG sentence classification assistant. Use only the provided English sentence translated from a Korean ESG disclosure. Return valid JSON.

\vspace{0.6em}
\textbf{[User]}\\
Task: Decide whether this corporate disclosure sentence discusses an ESG matter.

\vspace{0.3em}
Valid labels: ESG-related, Not ESG-related.

\vspace{0.3em}
Return JSON with this schema:\newline \texttt{\{"label": string\}}.

\vspace{0.3em}
Sentence: \texttt{\{sentence\}}

\vspace{0.6em}
\textbf{[Assistant]}\\
\texttt{\{"label": "\{label\}"\}}
\end{promptbox}

\subsection{ESG Topic Identification}
ESG topic identification is the task of classifying which ESG area or fine-grained issue an ESG-related text belongs to. This category includes ESG primary category classification, ESG issue classification, SASB issue topic classification, and TCFD recommendation classification. The model must identify the main topic of the input text and select the most appropriate label from predefined ESG categories or topic labels.
% ESG topic identification은 ESG 관련 텍스트가 어떤 ESG 영역 또는 세부 이슈에 해당하는지를 분류하는 task이다. 
% 이 category에는 ESG primary category classification, ESG issue classification, SASB issue topic classification, TCFD recommendation classification 등이 포함된다. 
% 모델은 입력 텍스트의 핵심 주제를 파악하고, 사전에 정의된 ESG category 또는 topic label 중 가장 적절한 항목을 선택해야 한다.

\begin{promptbox}{ESG Topic Identification}
\small
\textbf{[System]}\\
You are an ESG sentence classification assistant. Use only the provided English sentence translated from a Korean ESG disclosure. Return valid JSON.

\vspace{0.6em}
\textbf{[User]}\\
Task: Assign the most appropriate ESG pillar to the sentence.

\vspace{0.3em}
Allowed categories: Environmental, Social, Governance, Irrelevant.

\vspace{0.3em}
Return JSON with this schema:\newline \texttt{\{"category": string\}}.

\vspace{0.3em}
Sentence: \texttt{\{sentence\}}

\vspace{0.6em}
\textbf{[Assistant]}\\
\texttt{\{"category": "\{category\}"\}}
\end{promptbox}

\subsection{ESG Commitment Analysis}
ESG commitment analysis is the task of analyzing corporate ESG-related commitments, goals, action plans, reduction targets, and verifiability. This category includes climate commitment/action classification, net-zero reduction target classification, promise status classification, and promise verification timeline classification. The model must distinguish general ESG mentions from actual commitment expressions and assess their specificity, implementation status, target timeline, and verifiability.
% ESG commitment analysis는 기업의 ESG 관련 약속, 목표, 실행 계획, 감축 목표, 검증 가능성 등을 분석하는 task이다. 
% 이 category에는 climate commitment/action classification, net-zero reduction target classification, promise status classification, promise verification timeline classification 등이 포함된다. 
% 모델은 단순한 ESG 언급과 실제 commitment 표현을 구분하고, 해당 표현의 구체성, 이행 상태, 목표 시점, 검증 가능성 등을 판단해야 한다.

\begin{promptbox}{ESG Commitment Analysis}
\small
\textbf{[System]}\\
You are an ESG promise timeline classification assistant. Use only the provided English ESG report passage. Return valid JSON.

\vspace{0.6em}
\textbf{[User]}\\
Task: Classify the verification timeline for the ESG promise in the passage.

\vspace{0.3em}
Valid labels: Already, Less than 2 years, 2 to 5 years, More than 5 years, N/A.

\vspace{0.3em}
Use N/A when the passage does not contain an ESG promise, goal, target, commitment, or verifiable action.

\vspace{0.3em}
Return JSON with this schema:\\ \texttt{\{"verification\_timeline": string\}}.

\vspace{0.3em}
Passage: \texttt{\{passage\}}

\vspace{0.6em}
\textbf{[Assistant]}\\
\texttt{\{"verification\_timeline": "\{verification\_timeline\}"\}}
\end{promptbox}

\subsection{ESG Impact Analysis}
ESG impact analysis is the task of determining the direction, type, level, and duration of the impact that ESG-related events, activities, or statements have on companies, society, or the environment. This category includes climate sentiment classification, impact type classification, impact duration classification, and impact level classification. The model must identify the nature of ESG impact in the input text and analyze whether the impact is positive or negative, short-term or long-term, and how substantial it is.
% ESG impact analysis는 ESG 관련 사건, 활동, 진술이 기업, 사회, 환경에 미치는 영향의 방향, 유형, 수준, 지속 기간 등을 판단하는 task이다. 
% 이 category에는 climate sentiment classification, impact type classification, impact duration classification, impact level classification 등이 포함된다. 
% 모델은 입력 텍스트에서 ESG impact의 성격을 파악하고, 해당 영향이 긍정적 또는 부정적인지, 단기적 또는 장기적인지, 혹은 어느 정도 수준의 영향을 가지는지를 분석해야 한다.
\begin{promptbox}{ESG Impact Analysis}
\small
\textbf{[System]}\\
You are an ESG impact type classification assistant. Use only the provided English text. Return valid JSON.

\vspace{0.6em}
\textbf{[User]}\\
Task: Determine whether the text describes an ESG opportunity, risk, or cannot be distinguished.

\vspace{0.3em}
Allowed labels: Opportunity, Risk, Cannot distinguish.

\vspace{0.3em}
Return JSON with this schema:\newline \texttt{\{"impact\_type": string\}}.

\vspace{0.3em}
Title: \texttt{\{title\}}

\vspace{0.3em}
Text: \texttt{\{text\}}

\vspace{0.6em}
\textbf{[Assistant]}\\
\texttt{\{"impact\_type": "\{impact\_type\}"\}}
\end{promptbox}

\subsection{ESG Claim Verification}
ESG claim verification is the task of assessing the relationship between an ESG-related claim and its supporting evidence. This category includes climate claim verification, evidence quality classification, and evidence status classification. The model must evaluate whether a given claim is supported, refuted, difficult to determine, or sufficiently grounded by the evidence.
\begin{promptbox}{ESG Claim Verification}
\small
\textbf{[System]}\\
You are a climate claim verification assistant. Use only the provided claim and evidence. Do not use outside knowledge. Return valid JSON.

\vspace{0.6em}
\textbf{[User]}\\
Task: Determine the final verdict for the climate claim using only the provided evidence.

\vspace{0.3em}
Valid labels: SUPPORTS, REFUTES,\newline NOT\_ENOUGH\_INFO, DISPUTED.

\vspace{0.3em}
Return JSON with this schema:\newline \texttt{\{"label": string\}}.

\vspace{0.3em}
Claim: \texttt{\{claim\}}

\vspace{0.3em}
Evidence: \texttt{\{evidence\}}

\vspace{0.6em}
\textbf{[Assistant]}\\
\texttt{\{"label": "\{label\}"\}}
\end{promptbox}

\subsection{ESG Question Answering}
ESG question answering is the task of generating answers to ESG-related questions based on a given context, report, table, or evidence document. This category includes factoid QA, non-factoid QA, and table QA based on ESG reports and sustainability documents. The model must understand the intent of the question and generate an answer grounded in the provided context or evidence.

\begin{promptbox}{ESG Question Answering}
\small
\textbf{[System]}\\
You are a sustainability report question-answering assistant. Use only the provided context passage to answer descriptively. Return valid JSON.

\vspace{0.6em}
\textbf{[User]}\\
Task: Answer the ESG-related question based strictly on the provided context.

\vspace{0.3em}
Return JSON with this schema:\newline \texttt{\{"answer": string\}}.

\vspace{0.3em}
Context: \texttt{\{context\}}

\vspace{0.3em}
Question: \texttt{\{question\}}

\vspace{0.6em}
\textbf{[Assistant]}\\
\texttt{\{"answer": "\{answer\}"\}}
\end{promptbox}

\section{Implementation Details}
\label{app:implementation}
We compare models from three groups: large proprietary LLMs, open-source LLMs, and specialized ESG-oriented models. 
The proprietary LLM group includes recent API-based GPT and DeepSeek models. 
The open-source LLM group includes models from diverse families and sizes, such as Qwen, Llama, Gemma, Mistral, and DeepSeek-distilled models. 
The specialized model group includes SusGen-GPT. 
For \textsf{ESGPT}, we used Qwen3-8B as the base model and fine-tuned it using LoRA-based PEFT adapter training~\cite{hu2022lora}. 
During training, we used a maximum sequence length of 4096 tokens, 3 epochs, a per-device batch size of 1, gradient accumulation steps of 16, and an effective batch size of 16. 
The learning rate was set to $1\times10^{-4}$ with a cosine scheduler, a warmup ratio of 0.03, and no weight decay. 
Training was performed in bf16 precision with the AdamW optimizer. 
For LoRA, we used rank 64, alpha 128, dropout 0.05, and applied adapters to \texttt{q\_proj}, \texttt{k\_proj}, \texttt{v\_proj}, \texttt{o\_proj}, \texttt{gate\_proj}, \texttt{up\_proj}, and \texttt{down\_proj}. 
During inference, local models were loaded with 4-bit quantization by default and evaluated with a batch size of 4, a maximum input length of 4096 tokens, and maximum generation lengths of 256 tokens for classification tasks and 1024 tokens for QA or generation tasks. We used temperature 0.2 and top-$p$ 0.90 for decoding. API-based models were evaluated under the same temperature setting. All models were evaluated under the same evaluation protocol, with accuracy used as the evaluation metric across tasks. A JSON parser is first used to extract the required answer field from model outputs; if parsing fails, deterministic raw-text extraction is applied as a fallback. The extracted output is matched against the task-specific valid label set, and invalid or unmappable outputs are scored as incorrect. For ESG QA, formulated as multiple-choice QA, exact-match accuracy over predefined answer options is used. All experiments were conducted on a single NVIDIA A100 GPU. 

\section{Details of Experimental Results}
\label{app:additional_results}
This section provides additional detailed experimental results. While the main text focuses on comparisons with representative models on \textsf{ESGU-Bench}, this appendix reports the ranking of all models, category-level performance, and performance changes according to the training data scale in greater detail.
% 본 절에서는 세부 실험 결과를 추가로 제시한다. 
% 본문에서는 \textsf{ESGU-Bench}에서의 주요 모델들과의 비교 결과를 중심으로 논의하였으며, 본 부록에서는 전체 모델의 순위, category-level 성능, 그리고 학습 데이터 규모에 따른 성능 변화를 보다 상세히 보고한다. 

\subsection{\textsf{ESGU-Bench} Results}
\label{app:full_results}
Figure~\ref{fig:modelrank} shows the ranking of all 38 evaluated models on \textsf{ESGU-Bench} based on the overall score. The 38 models include the 35 LLMs discussed in the main text as well as additional \textsf{ESGPT} variants trained with different data scales, such as 10K, 20K and 50K. All \textsf{ESGPT} data-scale variants were trained with the same fixed random seed and training configuration to ensure a controlled comparison across data scales. In the overall ranking, the \textsf{ESGPT} models are placed among the top-performing models. Despite their 8B scale, they achieve higher overall scores than large proprietary LLMs, including DeepSeek-V4-Flash, DeepSeek-V4-Pro, GPT-4.1, and GPT-5.1. This suggests that, for ESG understanding tasks, instruction data designed for the ESG domain and task-specific adaptation can serve as important performance factors, beyond model size or general capability.

Table~\ref{tab:overall_performance} reports the full category-level performance across the six ESG understanding categories. For the final \textsf{ESGPT} result, we additionally report the mean $\pm$ standard deviation over five random seeds to assess performance stability. \textsf{ESGPT} achieves the highest overall score and records the best performance in ESG relevance detection, ESG topic identification, ESG commitment analysis, and ESG claim verification. This shows that \textsf{ESGPT} effectively adapts to tasks where understanding ESG disclosure contexts and task-specific label spaces is important, such as identifying ESG-related text, classifying ESG topics, interpreting corporate commitments and targets, and performing evidence-based verification. In contrast, DeepSeek-V4-Flash achieves the best performance in ESG impact analysis, while GPT-5.1 performs best in ESG question answering. This suggests that, for some tasks, the general reasoning ability and broad prior knowledge of large proprietary LLMs may still provide advantages.

Overall, \textsf{ESGPT}, despite being a relatively small 8B model, achieves performance that is competitive with or superior to that of large LLMs. This demonstrates that ESG-specialized models can be both efficient and effective. However, the overall results also show that both large proprietary LLMs and open-source LLMs still have limited performance in certain ESG categories. In particular, tasks that require complex contextual understanding and evidence-based judgment, such as ESG question answering, impact interpretation, and claim verification, leave substantial room for improvement. Therefore, \textsf{ESGU-Bench} can be viewed as a challenging benchmark that reveals both the limitations of current LLMs in ESG understanding tasks and directions for future improvement.

% Figure~\ref{fig:modelrank}는 \textsf{ESGU-Bench}에서 평가한 총 38개 모델의 overall score 기준 순위를 보여준다. 
% 전체 순위에서 \textsf{ESGPT} 계열 모델은 상위권에 위치하며, 8B 규모의 모델임에도 불구하고, DeepSeek-V4-Flash, DeepSeek-V4-Pro, GPT-4.1, GPT-5.1 등 대형 proprietary LLM보다 높은 overall score를 보였다. 
% 이는 ESG understanding task에서 모델 크기와 일반적인 성능보다 ESG 도메인에 맞게 구성된 instruction data와 task-specific adaptation이 중요한 성능 요인으로 작용할 수 있음을 보여준다.

% Table~\ref{tab:overall_performance}는 여섯 개 ESG understanding category에 대한 전체 category-level 성능을 제시한다. 
% \textsf{ESGPT}는 overall score에서 가장 높은 성능을 달성하였으며, ESG relevance detection, ESG topic identification, ESG commitment analysis, ESG claim verification에서 가장 높은 점수를 기록하였다. 
% 이는 ESG 관련 텍스트의 식별, 주제 분류, 기업의 약속 및 목표 해석, 근거 기반 검증과 같이 ESG 공시 문맥과 task-specific label space에 대한 이해가 중요한 과제에서 \textsf{ESGPT}가 효과적으로 적응했음을 보여준다. 
% 반면 ESG impact analysis에서는 DeepSeek-V4-Flash가, ESG reasoning에서는 GPT-5.1이 가장 높은 성능을 보였으며, 이는 일부 과제에서는 대형 proprietary LLM의 범용 추론 능력과 넓은 사전지식이 여전히 강점으로 작용할 수 있음을 시사한다.

% 종합적으로, \textsf{ESGPT}는 비교적 작은 8B 규모의 모델임에도 대형 LLM들과 경쟁하거나 이를 상회하는 성능을 달성하였다. 이는 ESG 특화 모델이 효율성과 효과성을 동시에 가질 수 있음을 보여준다. 
% 다만 전체 성능을 보면 대형 proprietary LLM과 open-source LLM 모두 특정 ESG category에서 여전히 제한적인 성능을 보이며, 특히 ESG reasoning, impact interpretation, claim verification과 같이 복합적인 문맥 이해와 근거 판단이 필요한 과제에서는 추가적인 성능 개선 여지가 크다. 따라서 \textsf{ESGU-Bench}는 현재 LLM들이 ESG 이해 과제에서 가지는 한계와 향후 개선 방향을 함께 드러내는 challenging benchmark로 볼 수 있다.

\subsection{Data Scale Analysis}
\label{app:scale_ablation}
Table~\ref{tab:scale_ablation} compares the performance of \textsf{ESGPT} under different training data scales using Qwen3-8B and Llama-3.1-8B-Instruct as base models. For both base models, the overall score generally improves as the training data scale increases from 10K to 30K and 50K. The Qwen3-8B-based model achieves the highest overall score of 67.00\% in the 50K setting, while the 30K setting shows a comparable score of 66.96\%. The Llama-3.1-8B-Instruct-based model also achieves its best performance in the 50K setting with 64.72\%, followed by the 30K setting with 64.24\%. In contrast, the full-data setting leads to a substantial drop in the overall score for both base models.

At the category level, the 30K or 50K settings show strong performance on most tasks, whereas ESG claim verification achieves the best performance in the 20K setting for both base models. This can be attributed to the relatively limited amount of training data available for ESG claim verification, which makes it difficult to maintain the balance of this task as the overall data scale increases. These results show that, rather than simply using more data, data construction and training strategies that consider task balance and label quality are critical for ESG instruction tuning.
% Table~\ref{tab:scale_ablation}는 Qwen3-8B와 Llama-3.1-8B-Instruct를 base model로 사용하여 학습 데이터 규모에 따른 \textsf{ESGPT}의 성능 변화를 비교한 결과를 보여준다. 두 base model 모두에서 학습 데이터가 10K에서 30K, 50K로 증가함에 따라 overall score가 전반적으로 향상되었다. Qwen3-8B 기반 모델은 50K 설정에서 67.00으로 가장 높은 overall score를 기록하였으며, 30K 설정도 66.96으로 유사한 수준의 성능을 보였다. Llama-3.1-8B-Instruct 기반 모델 역시 50K 설정에서 64.72로 가장 높은 성능을 보였고, 30K 설정이 64.24로 그 뒤를 이었다. 
% 반면 full-data 설정에서는 두 base model 모두 overall score가 크게 하락하였다.

% Category-level 결과, 대부분의 task에서는 30K 또는 50K 설정이 높은 성능을 보였으나, ESG claim verification은 두 base model 모두에서 20K 설정이 가장 높은 성능을 기록하였다. 이는 ESG claim verification에 활용 가능한 학습 데이터가 상대적으로 제한적이기 때문에, 전체 데이터 규모를 확대하는 과정에서 해당 task의 균형을 유지하기 어렵기 때문으로 해석할 수 있다. 즉, 단순히 더 많은 데이터를 사용하는 것보다 task balance, label quality를 고려한 데이터 구축 및 학습 방법의 효과를 확인하였다. 

\section{Qualitative Case Studies}
\label{app:case_studies}
This section presents five qualitative case studies observed in \textsf{ESGU-Bench}. Each case study includes the prompt, input text, gold label, and model responses used in the actual evaluation. Case Studies 1--3 are examples where \textsf{ESGPT} predicts the correct answer while large API-based LLMs fail, whereas Case Studies 4--5 are challenging cases where all comparison models fail. Through these cases, we analyze the strengths of ESG-specialized LLMs in interpreting ESG contexts, disclosure expressions, and task-specific label boundaries, as well as the types of ESG understanding that remain challenging for current models.

% 본 절에서는 \textsf{ESGU-Bench}에서 관찰된 5개의 qualitative case study를 제시한다. 
% 각 case study는 실제 평가에 사용된 prompt, 입력 텍스트, gold label, 그리고 모델 답변을 함께 보여준다. 
% Case Studies 1--3은 \textsf{ESGPT}가 정답을 예측한 반면 대형 API 기반 LLM들이 실패한 사례이며, Case Studies 4--5는 모든 비교 모델이 실패한 challenging cases이다. 
% 이를 통해 ESG 특화 LLM이 ESG 문맥, 공시 표현, task-specific label boundary를 해석하는 데 어떤 강점을 보이는지, 그리고 현재 모델들이 여전히 어려움을 겪는 ESG understanding 유형은 무엇인지 분석한다.

\begin{qualcasebox}{Case 1. Commitment Signal}
\small
\textbf{[User]}

\footnotesize
Classify whether this report excerpt includes an ESG target, commitment, or promise.\\
Allowed labels: Yes, No.

\vspace{0.6em}
\textbf{[Input]}

\footnotesize
Qualification, Orientation and Training All safety sensitive positions and supervisory staff receive extensive safety training, and all operations personnel are enrolled in a competency program that ensures only those deemed competent through training, exams, observation and sign-off by designated assessors can perform safety sensitive tasks unsupervised. Specific competency profiles are developed for each role.

\tcblower

\small
\textbf{[Gold and Model Answers]}

\vspace{0.2em}
\footnotesize
\begin{tabularx}{1.0\linewidth}{@{}Y C{0.25\linewidth} C{0.25\linewidth}@{}}
\hline
\textbf{Model} & \textbf{Answer} & \textbf{Result} \\
\hline
Gold & \goldans{Yes} & -- \\
\hline
DeepSeek-V4-Flash & \wrongans{No} & \xmarkred \\
DeepSeek-V4-Pro & \wrongans{No} & \xmarkred \\
GPT-4.1 & \wrongans{No} & \xmarkred \\
GPT-5.1 & \wrongans{No} & \xmarkred \\
\textsf{ESGPT} & \correctans{Yes} & \cmarkblue \\
\hline
\end{tabularx}
\end{qualcasebox}
Case 1 shows that ESG commitments are not always expressed through explicit goals, numerical targets, or declarative phrases such as ``we commit to.'' The input paragraph describes training, competency verification, testing, observation, and approval procedures for safety-sensitive roles, which can be viewed as ESG-related action signals indicating that the company operates a safety management system in practice. Large API-based LLMs answer No by focusing on the absence of explicit commitment expressions, whereas \textsf{ESGPT} answers Yes by capturing procedural and operational commitment signals that frequently appear in ESG reports. This demonstrates that an ESG-specialized LLM can more appropriately interpret implicit action signals in corporate disclosure texts.
% Case 1은 ESG commitment가 항상 명시적인 목표, 수치, 또는 ``we commit to''와 같은 선언형 표현으로만 나타나지 않는다는 점을 보여준다. 
% 입력 문단은 안전 민감 직무에 대한 교육, 역량 검증, 시험, 관찰, 승인 절차를 설명하며, 이는 기업이 안전 관리 체계를 실제로 운영하고 있음을 나타내는 ESG 관련 실행 신호로 볼 수 있다. 
% 대형 API 기반 LLM들은 명시적인 약속 표현이 없다는 점에 주목하여 No로 답변한 반면, \textsf{ESGPT}는 ESG 보고서에서 자주 나타나는 절차적,운영적 commitment signal을 포착하여 Yes로 답변하였다. 이는 ESG 특화 LLM이 기업 공시 문장에서 드러나는 암시적 실행 신호를 더 적절히 해석할 수 있음을 보여준다.


\begin{qualcasebox}{Case 2. Impact Severity}
\small
\textbf{[User]}

\footnotesize
Analyze the provided ESG news text and choose its impact severity.\\
Allowed labels: Low, Medium, High.

\vspace{0.6em}
\textbf{[Input]}

\footnotesize
\textbf{Title:} A.P. Møller Acquires Sustainable Food Packaging Company Faerch Group

\vspace{0.3em}
\textbf{Article:} Faerch is a leading European manufacturer of sustainable food packaging solutions within the Ready Meals, Protein, Food-to-Go and Dairy segments, providing circular solutions for rigid food packaging, with products made from up to 100\% recycled content and fully recyclable to new food grade products. According to A.P. Moller, Faerch is also the only food packaging company in the world to operate as an integrated recycler.

\tcblower

\small
\textbf{[Gold and Model Answers]}

\vspace{0.2em}
\footnotesize
\begin{tabularx}{1.0\linewidth}{@{}Y C{0.25\linewidth} C{0.25\linewidth}@{}}
\hline
\textbf{Model} & \textbf{Answer} & \textbf{Result} \\
\hline
Gold & \goldans{High} & -- \\
\hline
DeepSeek-V4-Flash & \wrongans{Medium} & \xmarkred \\
DeepSeek-V4-Pro & \wrongans{Medium} & \xmarkred \\
GPT-4.1 & \wrongans{Low} & \xmarkred \\
GPT-5.1 & \wrongans{Low} & \xmarkred \\
\textsf{ESGPT} & \correctans{High} & \cmarkblue \\
\hline
\end{tabularx}
\end{qualcasebox}
Case 2 shows that ESG impact analysis is not merely a task of determining ESG relevance, but also requires assessing the ESG significance and impact level of an event. The input article describes the acquisition of a sustainable food packaging company, which provides circular packaging solutions, products made from up to 100\% recycled content, full recyclability into food-grade packaging, and an integrated recycling system. API-based LLMs interpret this as a general ESG-related corporate acquisition and assign lower ratings of \texttt{Medium} or \texttt{Low}. In contrast, \textsf{ESGPT} predicts \texttt{High} by reflecting the strategic ESG significance of an event that combines circular economy practices with a recycling-based business transition. This suggests that ESG-specialized models may be better suited to assessing impact severity beyond the surface topic of an event.
% Case 2는 ESG impact analysis가 단순한 ESG 관련성 판단이 아니라, 사건의 ESG적 중요도와 파급 수준을 함께 평가해야 하는 과제임을 보여준다. 
% 입력 기사는 지속가능 식품 포장 기업의 인수 사례이며, 해당 기업은 순환형 포장 솔루션, 최대 100\% 재활용 소재, 식품 등급으로의 완전 재활용 가능성, 통합 재활용 시스템을 갖춘 것으로 설명된다. 
% API 기반 LLM들은 이를 ESG 관련성이 있는 일반적인 기업 인수로 해석하여 \texttt{Medium} 또는 \texttt{Low}로 낮게 평가하였다. 
% 반면 \textsf{ESGPT}는 순환경제와 재활용 기반 사업 전환이 결합된 사건의 전략적 ESG significance를 반영하여 \texttt{High}를 예측하였다. 
% 이는 ESG 특화 모델이 사건의 표면적 주제뿐 아니라 impact severity를 판단하는 데 더 적합할 수 있음을 시사한다.


\begin{qualcasebox}{Case 3. ESG Reasoning Boundary}
\small
\textbf{[User]}

\footnotesize
Answer the multiple-choice ESG reasoning question with a single option letter.\\
Valid options: A, B, C, D, Z.

\vspace{0.6em}
\textbf{[Input]}

\footnotesize
\textbf{Question:} What is the primary reason the ISSB rejected the hierarchical structure approach despite its potential to enhance comparability of digital information across sustainability-related disclosure standards, as described in the IFRS Taxonomy: IFRS Sustainability Disclosure Taxonomy 2024 -- IFRS S1 General Requirements for Disclosure of Sustainability-related Financial Information and IFRS S2 Climate-related Disclosures?

\vspace{0.3em}
\textbf{Options:}\\
A: The hierarchical structure would lead to fragmented tagged information that lacks sufficient context at lower levels.\\
B: Tagging systems are too new and untested to handle the complexity introduced by a hierarchical structure.\\
C: Regulators might require entities to use varying levels of elements, resulting in inconsistent digital information.\\
D: Elements at paragraph level were deemed insufficiently detailed, while lower-level elements lacked standalone utility.\\
Z: Not sure

\tcblower

\small
\textbf{[Gold and Model Answers]}

\vspace{0.2em}
\footnotesize
\begin{tabularx}{1.0\linewidth}{@{}Y C{0.25\linewidth} C{0.25\linewidth}@{}}
\hline
\textbf{Model} & \textbf{Answer} & \textbf{Result} \\
\hline
Gold & \goldans{C} & -- \\
\hline
DeepSeek-V4-Flash & \wrongans{A} & \xmarkred \\
DeepSeek-V4-Pro & \wrongans{A} & \xmarkred \\
GPT-4.1 & \wrongans{A} & \xmarkred \\
GPT-5.1 & \wrongans{A} & \xmarkred \\
\textsf{ESGPT} & \correctans{C} & \cmarkblue \\
\hline
\end{tabularx}
\end{qualcasebox}
Case 3 shows that, in ESG question answering, it is important to distinguish between explanations that are generally plausible and those that align with the correct answer criterion. Option A, selected by the API models, is a plausible explanation in that hierarchical tagging may cause context loss at lower-level tags. However, the correct answer required by this question is not about the general limitation of hierarchical structures, but about the reduced consistency of digital information when regulators require the use of taxonomy elements at different levels. \textsf{ESGPT} predicts C, which aligns with the benchmark's answer criterion, rather than the generally plausible answer selected by the API models. This shows that an ESG-specialized LLM can better distinguish subtle semantic differences among options and task-specific decision boundaries in ESG question answering tasks.

% Case 3은 ESG reasoning에서 일반적으로 그럴듯한 설명과 정답 기준에 부합하는 설명을 구분하는 능력이 중요함을 보여준다. 
% API 모델들이 선택한 A는 hierarchical tagging이 낮은 수준의 태그에서 문맥을 잃게 만들 수 있다는 점에서 자연스러운 설명이다. 
% 그러나 이 문항에서 요구하는 정답은 hierarchical structure의 일반적 한계가 아니라, 규제기관이 서로 다른 수준의 taxonomy element 사용을 요구할 경우 digital information의 일관성이 낮아질 수 있다는 점이다. 
% \textsf{ESGPT}는 API 모델들이 선택한 일반적으로 타당해 보이는 답변이 아니라, benchmark의 정답 기준에 부합하는 C를 예측하였다. 
% 이는 ESG 특화 LLM이 ESG reasoning task에서 선택지 간 미세한 의미 차이와 task-specific decision boundary를 더 잘 구분할 수 있음을 보여준다.

\begin{qualcasebox}{Case 4. Commitment Boundary}
\small
\textbf{[User]}

\footnotesize
Review the paragraph and return whether it concerns climate commitments and actions.\\
Valid labels: yes, no.

\vspace{0.6em}
\textbf{[Input]}

\footnotesize
Furthermore, according to the United Nations Environment Programme's (UNEP) latest Emissions Gap Report, 1.5$^\circ$C will involve ``expanding renewable energy for electrification, phasing out coal for rapid decarbonisation of the energy system, decarbonising transport with a focus on electric mobility, decarbonising energy-intensive industry and avoiding future emissions, while improving energy access''; all very challenging actions.

\tcblower

\small
\textbf{[Gold and Model Answers]}

\vspace{0.2em}
\footnotesize
\begin{tabularx}{1.0\linewidth}{@{}Y C{0.25\linewidth} C{0.25\linewidth}@{}}
\hline
\textbf{Model} & \textbf{Answer} & \textbf{Result} \\
\hline
Gold & \goldans{no} & -- \\
\hline
DeepSeek-V4-Flash & \wrongans{yes} & \xmarkred \\
DeepSeek-V4-Pro & \wrongans{yes} & \xmarkred \\
GPT-4.1 & \wrongans{yes} & \xmarkred \\
GPT-5.1 & \wrongans{yes} & \xmarkred \\
\textsf{ESGPT} & \wrongans{yes} & \xmarkred \\
\hline
\end{tabularx}
\end{qualcasebox}
Case 4 shows a case where all models overreact to climate-action expressions. The input paragraph contains strong climate-action terms such as renewable energy, phasing out coal, decarbonising transport, and electric mobility. However, the sentence does not describe a commitment by a specific company or reporting entity, but rather a general climate transition pathway presented in a UNEP report. In other words, it explains the need for climate action, but does not indicate an explicit commitment or action by the reporting entity. This case shows that distinguishing salient climate keywords from the actual commitment boundary remains challenging for all models, including \textsf{ESGPT}.

% Case 4는 모든 모델이 기후 행동 표현에 과하게 반응한 경우이다. 
% 입력 문단에는 renewable energy, phasing out coal, decarbonising transport, electric mobility와 같은 강한 climate-action 표현이 포함되어 있다. 
% 그러나 이 문장은 특정 기업이나 보고 주체가 수행하겠다고 밝힌 약속이 아니라, UNEP 보고서가 제시하는 일반적인 기후 전환 경로를 설명한다. 
% 즉, 기후 행동의 필요성을 설명하는 문장이지, reporting entity의 명시적 commitment/action을 나타내는 문장은 아니다. 
% 이 사례는 주요 키워드와 실제 commitment boundary를 구분하는 것이 \textsf{ESGPT}를 포함한 모든 모델에게 여전히 어려운 문제임을 보여준다.

\begin{qualcasebox}{Case 5. Sentiment Boundary}
\small
\textbf{[User]}

\footnotesize
Classify the climate sentiment of the given paragraph.\\  
Allowed labels: risk, neutral, opportunity.

\vspace{0.6em}
\textbf{[Input]}

\footnotesize
The important initiatives taken by DNP Group in moving toward a decarbonized society include reducing the amount of CO2-generating energy used, namely energy conservation, switching to low CO2-emission fuels, and introducing renewable energy sources.

\tcblower

\small
\textbf{[Gold and Model Answers]}

\vspace{0.2em}
\footnotesize
\begin{tabularx}{1.0\linewidth}{@{}Y C{0.25\linewidth} C{0.25\linewidth}@{}}
\hline
\textbf{Model} & \textbf{Answer} & \textbf{Result} \\
\hline
Gold & \goldans{neutral} & -- \\
\hline
DeepSeek-V4-Flash & \wrongans{opportunity} & \xmarkred \\
DeepSeek-V4-Pro & \wrongans{opportunity} & \xmarkred \\
GPT-4.1 & \wrongans{opportunity} & \xmarkred \\
GPT-5.1 & \wrongans{opportunity} & \xmarkred \\
\textsf{ESGPT} & \wrongans{opportunity} & \xmarkred \\
\hline
\end{tabularx}
\end{qualcasebox}
Case 5 shows a case where all models interpret positive climate-action expressions as \texttt{opportunity}. The input paragraph describes how DNP Group implemented energy savings, a shift to low-carbon fuels, and the adoption of renewable energy to move toward a decarbonized society. However, the paragraph does not explicitly describe business opportunities such as corporate growth, cost reduction, competitiveness enhancement, or market expansion. Therefore, under the benchmark criterion, the sentence is closer to a neutral description of climate-related actions, and the gold label is \texttt{neutral}. This case shows that distinguishing mentions of positive ESG actions from opportunity framing from a corporate perspective remains challenging.

Overall, Cases 1--3 show that ESG-specialized LLMs can more appropriately interpret implicit signals in ESG contexts, the impact severity of events, and subtle semantic differences among answer options. In contrast, Cases 4--5 show that current models may still overreact to ESG/climate-related keywords or positive expressions, and may struggle to distinguish explicit commitments from general climate-action statements or to identify sentiment label boundaries. These findings confirm that \textsf{ESGU-Bench} provides an evaluation setting that requires fine-grained semantic distinctions in ESG contexts and task-specific judgment, beyond simple ESG keyword detection. At the same time, current models still show limitations in fine-grained ESG judgments, such as explicit commitment detection, sentiment framing, and reasoning boundaries, suggesting the need to improve ESG-specific reasoning and label-boundary awareness.

% Case 5는 모든 모델이 긍정적인 기후 행동 표현을 \texttt{opportunity}로 해석한 경우이다. 
% 입력 문단은 DNP Group이 탈탄소 사회로 나아가기 위해 에너지 절약, 저탄소 연료 전환, 재생에너지 도입 등을 수행했다는 내용을 설명한다. 
% 그러나 해당 문장은 기업의 성장 기회, 비용 절감, 경쟁력 강화, 시장 확대와 같은 business opportunity를 명시적으로 서술하지 않는다. 
% 따라서 benchmark 기준에서는 기후 관련 조치를 중립적으로 설명하는 문장에 가깝고, gold label은 \texttt{neutral}이다. 
% 이 사례는 긍정적인 ESG 행동의 언급과 기업 관점의 opportunity framing을 구분하는 것이 여전히 어려운 지점임을 보여준다.

% 종합적으로, Case 1,2,3은 ESG 특화 LLM이 ESG 문맥에서 나타나는 암시적 신호, 사건의 impact severity, 선택지 간 세부 의미 차이를 더 적절히 해석할 수 있음을 보여준다. 
% 반면 Case 4와 5는 현재 모델들이 여전히 ESG/기후 관련 키워드나 긍정적 표현에 과하게 반응할 수 있으며, 명시적 commitment 여부와 sentiment label boundary를 구분하는 데 어려움을 겪을 수 있음을 보여준다. 따라서 \textsf{ESGU-Bench}는 단순한 ESG 키워드 탐지를 넘어, ESG 문맥의 세부적인 의미 구분과 task-specific 판단 능력을 요구하는 평가 환경임을 확인할 수 있다. 동시에 현재 모델들은 명시적 commitment 여부, sentiment framing, reasoning boundary와 같은 세밀한 ESG 판단에서 여전히 한계를 보이며, 이는 향후 ESG-specific reasoning 및 label boundary 인식 능력 개선이 필요함을 시사한다.

\section{TCFD Report Generation}
\label{app:report_generation}
We qualitatively compare the TCFD-style reports generated by \textsf{ESGPT} and SusGen-GPT at the report level. A TCFD-style report is constructed by generating answers to 11 questions corresponding to the major disclosure areas of the TCFD recommendations, using evidence retrieved from corporate sustainability/ESG reports, and organizing them into a single report. SusGen-GPT is an LLM specialized for sustainability report generation, trained with both financial NLP tasks and ESG report generation, and is used as a strong domain-specific baseline~\cite{wu2025susgen}.

For analysis, we use NYSE\_PBR\_2016, which corresponds to the 2016 report of Petrobras, a company listed on the New York Stock Exchange~\cite{ni2023chatreport}. This report contains ESG information related to governance, compliance, and risk management, but climate-specific disclosure is unclear for some TCFD items. It is therefore suitable for examining how the two models use general ESG evidence when generating climate-related disclosure. Table~\ref{tab:pbr_side_by_side_report} compares the generation results of the two models on NYSE\_PBR\_2016.

The results show that \textsf{ESGPT} tends to preserve the limitations of the provided evidence when climate-specific disclosure is unclear. For example, in items such as board oversight, management roles, and risk identification processes, \textsf{ESGPT} does not directly infer climate-related oversight or climate-specific risk processes from general governance or risk management information. In contrast, SusGen-GPT generates longer and more report-like responses, but sometimes extends general ESG information into climate-specific disclosure.

This difference is particularly evident in the TCFD 10 item on Scope emissions. \textsf{ESGPT} states that Scope 1 and Scope 2 emissions are disclosed, while distinguishing that Scope 3 emissions are not explicitly mentioned. In contrast, SusGen-GPT first states that the organization discloses Scope 1, Scope 2, and Scope 3 emissions, and then adds within the same answer that explicit Scope 3 data are not provided. This overstates disclosure coverage and introduces an internal inconsistency. Similar tendencies are observed in TCFD 1 and TCFD 9: SusGen-GPT links the Petrobras Corruption Prevention Program to climate-related risk oversight and states that a comprehensive set of metrics exists despite the lack of specific quantitative indicators. By contrast, \textsf{ESGPT} uses more constrained wording when the evidence is insufficient to support a climate-specific claim.

Overall, SusGen-GPT has the advantage of producing report-style responses, but it may over-claim or generate internally inconsistent statements when the evidence is indirect or incomplete. In contrast, \textsf{ESGPT} more carefully preserves ESG disclosure boundaries and produces more evidence-grounded, conservative reports. Since ESG report generation is closely tied to corporate disclosure, these results suggest that faithfulness to evidence is critical for reliable ESG report generation.

% \textsf{ESGPT}와 SusGen-GPT가 생성한 TCFD-style report를 보고서 수준에서 비교하여, 두 모델의 ESG 보고서 생성 양상 차이를 정성적으로 분석한다. TCFD-style report는 기업의 지속가능성/ESG 보고서에서 검색된 evidence를 바탕으로, TCFD 권고안의 주요 disclosure 영역에 대응하는 11개 질문에 대한 답변을 생성하고 이를 하나의 보고서 형태로 구성한 것이다. SusGen-GPT는 금융 NLP task와 ESG report generation을 함께 고려하여 학습된 지속가능성 보고서 생성 특화 LLM으로, 본 비교에서 강한 도메인 특화 baseline으로 사용된다~\cite{wu2025susgen}.

% 분석 대상으로는 뉴욕증권거래소 상장 기업인 Petrobras의 2016년 보고서에 해당하는 NYSE\_PBR\_2016을 사용하였다~\cite{ni2023chatreport}. 이 보고서는 governance, compliance, risk management 관련 ESG 정보는 포함하지만, 일부 TCFD 항목에서는 climate-specific disclosure가 명확하지 않다. 따라서 제한적인 evidence 조건에서 두 모델이 일반 ESG 정보를 어떻게 보고서 문장으로 구성하고, 이를 climate-specific disclosure로 확장하는 방식에 어떤 차이를 보이는지 확인하기에 적합하다. Table~\ref{tab:pbr_side_by_side_report}에서 NYSE\_PBR\_2016에 대한 두 모델의 생성 결과를 비교한다.

% 생성 결과, \textsf{ESGPT}는 climate-specific disclosure가 명확하지 않은 항목에서 제공된 evidence의 한계를 비교적 일관되게 유지하는 경향을 보인다. 예를 들어 이사회 감독, 경영진 역할, 위험 식별 프로세스와 같은 항목에서 \textsf{ESGPT}는 일반적인 governance 또는 risk management 정보가 존재하더라도, 이를 곧바로 climate-related oversight나 climate-specific risk process로 단정하지 않는다. 반면 SusGen-GPT는 더 긴 답변을 생성하지만, 일부 항목에서는 일반 ESG 정보를 climate-specific disclosure로 확장하는 경향이 나타난다.

% 이러한 차이는 특히 TCFD 10의 Scope emissions 항목에서도 드러난다. \textsf{ESGPT}는 Scope 1과 Scope 2 배출량이 공개되었음을 설명하면서도, Scope 3 배출량은 명시적으로 언급되지 않았다고 구분한다. 반면 SusGen-GPT는 조직이 Scope 1, Scope 2, Scope 3 배출량을 모두 공개한다고 먼저 서술한 뒤, 같은 답변 안에서 명시적인 Scope 3 데이터는 제공되지 않는다고 덧붙인다. 
% 이는 disclosure coverage를 과도하게 인정하는 동시에, 답변 내부에서 서로 충돌하는 진술이 발생한 사례로 볼 수 있다.
% 또한 TCFD 1과 TCFD 9에서도 유사한 경향이 관찰된다. SusGen-GPT는 Petrobras Corruption Prevention Program을 climate-related risk oversight와 연결하거나, 구체적인 정량 지표가 부족함에도 comprehensive set of metrics가 있다고 서술한다. 이와 달리 \textsf{ESGPT}는 해당 정보가 climate-specific claim을 뒷받침하기에 충분하지 않은 경우 제한적인 표현을 사용한다.

% 종합하면, SusGen-GPT는 보고서 다운 문장을 생성하는 장점이 있지만, evidence가 간접적이거나 불완전한 상황에서는 일반 ESG 정보를 climate-specific disclosure로 과도하게 확장하거나, 답변 내부에서 서로 충돌하는 진술을 생성할 수 있다. 
% 반면 \textsf{ESGPT}는 제공된 evidence의 범위를 신중하게 유지하며, 명시적으로 확인되지 않은 정보에 대해서는 제한적인 표현을 사용한다. 
% ESG 보고서 생성은 기업 공시와 연결되는 작업이므로, 유창성보다 사실 기반의 정확성과 근거 충실성이 중요하다. 
% 이러한 맥락에서 over-claiming이나 internal inconsistency는 보고서 신뢰성을 저해하는 치명적인 hallucination 문제가 될 수 있다. 
% 본 사례는 \textsf{ESGPT}가 ESG disclosure boundary를 더 잘 구분하고, 근거 중심의 보수적인 보고서 생성을 수행함으로써 보다 신뢰 가능한 ESG report generation에 적합함을 보여준다.
\begin{table*}[t]
\centering
\scriptsize
\setlength{\tabcolsep}{4pt}
\renewcommand{\arraystretch}{1.12}
\resizebox{\textwidth}{!}{%
\begin{tabular}{p{0.48\textwidth}|p{0.48\textwidth}}
\hline
\textbf{\textsf{ESGPT}} & \textbf{SusGen-GPT} \\
\hline

\multicolumn{2}{l}{\textbf{TCFD 1. Board Oversight of Climate-related Risks and Opportunities}} \\
\hline
The board oversees climate-related risks and opportunities indirectly through its focus on operational safety and financial leverage reduction. 
\textcolor{caseblue}{However, the sustainability report does not explicitly link these strategies to climate-specific measures or provide details on how the board directly addresses climate risks.}
&
The board oversees climate-related risks and opportunities through strategic goals and governance mechanisms. 
It also monitors the Petrobras Corruption Prevention Program, which \textcolor{casered}{likely includes climate-related risks}. 
\textcolor{casered}{However, the evidence does not explicitly state that climate-related risks are addressed in those meetings or reports.}
\\
\hline

\multicolumn{2}{l}{\textbf{TCFD 2. Management Role in Assessing and Managing Climate-related Risks}} \\
\hline
\textcolor{caseblue}{The evidence does not explicitly link management's role to climate-related risks and opportunities.} 
The disclosure focuses instead on general business strategy, management system outcomes, compliance, governance, prospects, and challenges.
&
Management plays a significant role in assessing and managing climate-related risks and opportunities. 
This is inferred from sections discussing compliance, prospects, challenges, governance, and management system results. 
However, the report does not provide detailed strategies or quantitative evidence for how these risks are addressed.
\\
\hline

\multicolumn{2}{l}{\textbf{TCFD 3. Climate-related Risks and Opportunities over Time Horizons}} \\
\hline
The organisation identifies risks that may jeopardize strategic goals, business execution, regulatory compliance, ethics, market conditions, credit, and liquidity. 
These risks are prioritized based on potential consequences and probability of occurrence.
&
The organization identifies climate-related risks across short-, medium-, and long-term horizons. 
Short-term risks include operational risks, medium-term risks involve compliance, and long-term risks relate to strategic goals. 
\textcolor{casered}{However, explicit climate-related opportunities are not identified in the evidence.}
\\
\hline

\multicolumn{2}{l}{\textbf{TCFD 4. Impact on Strategy, Financial Performance, and Planning}} \\
\hline
Climate-related risks and opportunities are discussed through broader operational challenges, regulatory compliance, and market volatility. 
\textcolor{caseblue}{The evidence is vague and focuses more on general compliance and operational risks than on specific climate-related financial impacts.}
&
\textcolor{casered}{Climate-related risks and opportunities significantly impact} the organization's business strategy, economic and financial performance, and financial planning. 
The answer links this claim to general risks and management-system disclosures, but notes that quantitative financial evidence is not provided.
\\
\hline

\multicolumn{2}{l}{\textbf{TCFD 5. Strategy Resilience under Climate Scenarios}} \\
\hline
The organisation's strategy is aligned with its risk management system and integrated into strategic decisions. 
\textcolor{caseblue}{However, the evidence does not explicitly address resilience under different climate scenarios, including a 2$^\circ$C or lower scenario, or physical climate risks.}
&
The strategy shows integrated risk management and alignment with the Strategic Plan. 
However, the evidence does not provide specific details on climate physical risks or resilience under a 2$^\circ$C scenario.
\\
\hline

\multicolumn{2}{l}{\textbf{TCFD 6. Processes for Identifying and Assessing Climate-related Risks}} \\
\hline
The organisation identifies risks through interviews, online research, and strategic analysis. 
Risks are prioritized based on consequences and probability of occurrence. 
\textcolor{caseblue}{The process for identifying climate-related risks is not explicitly detailed.}
&
The organization uses a multi-step process involving broad diagnosis, online research, face-to-face interviews with executives, Strategic Plan alignment, business-plan risk identification, social-risk analysis, and board monitoring. 
\textcolor{casered}{However, the evidence does not provide specific details on how climate-related risks are identified within the broader risk process.}
\\
\hline

\multicolumn{2}{l}{\textbf{TCFD 7. Management of Climate-related Risks}} \\
\hline
The organization manages climate-related risks by monitoring global alternatives for low-carbon energy, adapting projects and facilities, and disclosing annual atmospheric emissions data from nearly 30,000 sources. 
These efforts were recognized through a high CDP score in 2016.
&
The organization manages climate-related risks through strategic planning, risk responses, monitoring low-carbon energy models, and adapting projects and facilities. 
The answer also notes that the evidence does not provide detailed quantitative strategies for climate-risk mitigation.
\\
\hline

\multicolumn{2}{l}{\textbf{TCFD 8. Integration into Overall Risk Management}} \\
\hline
Climate-related risks are integrated into overall risk management through alignment with the Strategic Plan, risk identification and prioritization, and structured tools such as a social risk matrix. 
These processes support systematic consideration of risks in strategic decisions.
&
Climate-related risks are integrated into overall risk management through Strategic Plan alignment, periodic reporting to the Executive Committee of Risks, social-risk structures, and executive risk oversight. 
\textcolor{casered}{However, the evidence does not explicitly mention climate-related risk integration, only broader risk integration.}
\\
\hline

\multicolumn{2}{l}{\textbf{TCFD 9. Metrics for Assessing Climate-related Risks and Opportunities}} \\
\hline
\textcolor{caseblue}{The organisation does not explicitly disclose specific metrics} for assessing climate-related risks and opportunities. 
Instead, it focuses on broader governance and compliance strategies, which limits clarity on how climate risks are quantified or managed.
&
The organization uses \textcolor{casered}{a comprehensive set of metrics} to assess climate-related risks and opportunities. 
However, \textcolor{casered}{the specific metrics and quantitative values are not explicitly detailed} in the evidence, limiting a thorough assessment.
\\
\hline

\multicolumn{2}{l}{\textbf{TCFD 10. Scope 1, Scope 2, and Scope 3 GHG Emissions}} \\
\hline
The organisation discloses Scope 1 and Scope 2 emissions. 
Scope 1 emissions decreased from 77.4 million metric tons in 2015 to 66.1 million metric tons in 2016, while Scope 2 emissions decreased from 0.7 to 0.4 million metric tons. 
\textcolor{caseblue}{Scope 3 emissions are not explicitly mentioned.}
&
The organization \textcolor{casered}{discloses Scope 1, Scope 2, and Scope 3 greenhouse gas emissions}. 
Scope 1 emissions were 66.1 million tCO2e and Scope 2 emissions were 0.4 million tCO2e in 2016. 
\textcolor{casered}{The report does not provide explicit Scope 3 emissions data.}
\\
\hline

\multicolumn{2}{l}{\textbf{TCFD 11. Targets and Performance against Climate-related Risks and Opportunities}} \\
\hline
\textcolor{caseblue}{The organization does not explicitly mention specific targets} for understanding, quantifying, or benchmarking climate-related risks and opportunities. 
Instead, it focuses on broader sustainability initiatives and governance or compliance strategies.
&
The organization provides governance, business strategy, and management system information. 
However, the evidence does not explicitly state specific climate-related risks and opportunities targets or KPIs, making it unclear how the organization performs against them.
\\
\hline

\end{tabular}%
}
\caption{Side-by-side comparison of TCFD-style report generation outputs from \textsf{ESGPT} and SusGen-GPT for NYSE\_PBR\_2016. The outputs are compressed for readability, with blue highlights indicating evidence-bounded statements and red highlights indicating potential hallucination risks.}
\label{tab:pbr_side_by_side_report}
\end{table*}

%------------------------------------
\begin{figure*}[t]
    \centering
    \includegraphics[width=1\linewidth]{Figure/APFig4_Modelrank.pdf}
    \caption{Overall model ranking on \textsf{ESGU-Bench} across 38 evaluated models. Bars indicate overall scores, and colors denote model groups: proprietary LLMs, specialized LLMs, and open-source LLMs.}
    \label{fig:modelrank}
\end{figure*}
%------------------------------------
\begin{table*}[t]
\centering
\scriptsize
\setlength{\tabcolsep}{3.0pt}
\renewcommand{\arraystretch}{1.08}
\resizebox{\textwidth}{!}{%
\begin{tabular}{l|l|c|cccccc|c}
\hline
\textbf{Type} & \textbf{Model} & \textbf{Size}
& \makecell[c]{\textbf{ESG Relevance}\\\textbf{Detection}}
& \makecell[c]{\textbf{ESG Topic}\\\textbf{Identification}}
& \makecell[c]{\textbf{ESG Commitment}\\\textbf{Analysis}}
& \makecell[c]{\textbf{ESG Impact}\\\textbf{Analysis}}
& \makecell[c]{\textbf{ESG Claim}\\\textbf{Verification}}
& \makecell[c]{\textbf{ESG}\\\textbf{QA}}
& \textbf{Overall} \\
\hline

\multirow{28}{*}{\makecell[c]{Open\\Source\\LLM}}
& Qwen2.5 & 0.5B & 0.00 & 0.04 & 0.06 & 0.00 & 0.17 & 53.87 & 9.02 \\
& Qwen2.5 & 1.5B & 0.19 & 0.13 & 0.06 & 0.00 & 1.80 & 54.84 & 9.50 \\
& Qwen2.5 & 3B & 55.43 & 22.26 & 53.03 & 57.90 & 44.37 & 58.71 & 48.62 \\
& Qwen2.5 & 7B & 64.59 & 35.31 & 55.66 & 59.63 & 49.68 & 62.76 & 54.61 \\
& Qwen2.5 & 14B & 69.78 & 32.93 & 61.32 & 63.33 & 52.55 & 61.97 & 56.98 \\
& Qwen2.5 & 32B & 80.26 & 36.36 & 57.90 & \underline{68.27} & 52.98 & 60.39 & 59.36 \\
& Qwen2.5 & 72B & 76.77 & 37.49 & 62.33 & 68.02 & 55.89 & 61.88 & 60.40 \\

& Qwen3 & 0.6B & 64.12 & 26.11 & 49.50 & 59.51 & 40.39 & 13.20 & 42.14 \\
& Qwen3 & 1.7B & 60.06 & 38.41 & 60.48 & 58.02 & 50.45 & 47.62 & 52.51 \\
& Qwen3 & 4B & 67.14 & 42.89 & 62.50 & 67.65 & 53.06 & 61.88 & 59.19 \\
& Qwen3 & 8B & 73.28 & 43.61 & 65.70 & 59.14 & 53.64 & 60.21 & 59.26 \\

& Llama-3.1 & 8B & 75.26 & 37.28 & 67.00 & 56.66 & 53.21 & 62.50 & 58.65 \\
& Llama-3.2 & 1B & 49.76 & 14.69 & 42.77 & 49.75 & 37.90 & 54.84 & 41.62 \\
& Llama-3.2 & 3B & 66.76 & 34.98 & 60.20 & 56.91 & 48.22 & 57.83 & 54.15 \\
& Llama-3.3 & 70B & \underline{81.87} & 38.16 & 63.62 & 66.42 & 55.50 & 65.58 & 61.86 \\

& Gemma-3 & 1B & 45.14 & 18.54 & 41.14 & 64.57 & 38.24 & 29.05 & 39.45 \\
& Gemma-3 & 4B & 61.95 & 28.16 & 57.34 & 62.47 & 40.17 & 48.42 & 49.75 \\
& Gemma-3 & 12B & 72.62 & 34.64 & 66.26 & 67.41 & 56.15 & 24.21 & 53.55 \\
& Gemma-3 & 27B & 80.36 & 35.23 & 67.10 & 65.56 & 58.93 & 45.95 & 58.85 \\

& Mistral-v0.1 & 7B & 69.69 & 25.56 & 54.32 & 57.16 & 47.07 & 58.98 & 52.13 \\
& Mistral-v0.2 & 7B & 69.41 & 32.47 & 57.51 & 64.07 & 50.06 & 59.77 & 55.55 \\
& Mistral-v0.3 & 7B & 72.33 & 29.12 & 56.67 & 63.33 & 51.61 & 60.04 & 55.52 \\

& DeepSeek-R1-Distill & 1.5B & 21.25 & 6.69 & 17.32 & 22.47 & 3.25 & 4.58 & 12.59 \\
& DeepSeek-R1-Distill & 7B & 27.38 & 7.15 & 22.25 & 13.21 & 5.05 & 55.37 & 21.74 \\
& DeepSeek-R1-Distill & 14B & 47.88 & 12.13 & 32.96 & 22.84 & 9.42 & 62.85 & 31.35 \\
& DeepSeek-R1-Distill & 32B & 35.03 & 10.00 & 26.46 & 21.73 & 7.07 & 63.12 & 27.23 \\
& DeepSeek-R1-Distill & 8B & 48.73 & 12.72 & 28.53 & 20.25 & 7.41 & 57.31 & 29.16 \\
& DeepSeek-R1-Distill & 70B & 43.63 & 11.38 & 27.13 & 14.32 & 4.71 & 65.05 & 27.70 \\
\hline

\multirow{4}{*}{\makecell[c]{Proprietary\\LLM}}
& DeepSeek-V4-Flash & 284B & 82.63 & 48.87 & 66.26 & \textbf{71.11} & \underline{58.42} & 67.69 & \underline{65.83} \\
& DeepSeek-V4-Pro & 1600B & 79.70 & 46.82 & 66.26 & 66.79 & 56.32 & 68.05 & 63.99 \\
& GPT-4.1 & 2200B$^{\dagger}$ & 81.30 & 46.82 & \underline{68.67} & 65.06 & 47.79 & \underline{69.19} & 63.14 \\
& GPT-5.1 & 1300B$^{\dagger}$ & 79.13 & \underline{49.92} & 67.66 & 62.59 & 45.52 & \textbf{72.89} & 62.95 \\
\hline

\multirow{3}{*}{\makecell[c]{Specialized\\LLM}}
& SusGen-GPT-Llama & 8B & 67.61 & 37.70 & 62.00 & 53.70 & 56.66 & 57.66 & 55.89 \\
& SusGen-GPT-Mistral & 7B & 71.20 & 34.10 & 58.52 & 62.84 & 52.98 & 61.44 & 56.85 \\
& \textsf{ESGPT}~(Ours) & 8B & \textbf{84.32~$\pm$~(0.09)} & \textbf{54.94~$\pm$~(0.16)} & \textbf{73.37~$\pm$~(0.31)} & 65.93~$\pm$~(0.17) & \textbf{59.40~$\pm$~(0.05)} & 63.82~$\pm$~(0.18) & \textbf{66.96~$\pm$~(0.10)}$^{*}$ \\
\hline

\end{tabular}%
}
\caption{Overall performance comparison on \textsf{ESGU-Bench}. The results are grouped by model type, with task-level scores reported across six ESG understanding categories. Best scores are shown in bold, second-best scores are underlined, and model sizes marked with $^{\dagger}$ indicate estimated parameter counts. To assess the performance stability of \textsf{ESGPT}, we report the mean $\pm$ standard deviation over five random seeds after excluding the highest and lowest seed results. The $^{*}$ marker in the Overall column denotes t-test significance against the second-best model, DeepSeek-V4-Flash, with $p < 0.001$.}
\label{tab:overall_performance}
\end{table*}
%------------------------------------
\begin{table*}[t]
\centering
\scriptsize
\setlength{\tabcolsep}{3.0pt}
\renewcommand{\arraystretch}{1.08}
\resizebox{\textwidth}{!}{%
\begin{tabular}{l|l|c|cccccc|c}
\hline
\textbf{Base Model} 
& \textbf{Model} 
& \textbf{Train Time}
& \makecell[c]{\textbf{ESG Relevance}\\\textbf{Detection}}
& \makecell[c]{\textbf{ESG Topic}\\\textbf{Identification}}
& \makecell[c]{\textbf{ESG Commitment}\\\textbf{Analysis}}
& \makecell[c]{\textbf{ESG Impact}\\\textbf{Analysis}}
& \makecell[c]{\textbf{ESG Claim}\\\textbf{Verification}}
& \makecell[c]{\textbf{ESG}\\\textbf{QA}}
& \textbf{Overall} \\
\hline

\multirow{5}{*}{\makecell[c]{Qwen3-8B}}
& ESGPT-10K & 03h 05m 49s & 81.68 & 53.72 & 69.28 & 63.58 & \underline{60.34} & 62.15 & 65.13 \\
& ESGPT-20K & 06h 41m 06s & 82.53 & 53.93 & 69.62 & 63.95 & \textbf{60.39} & 63.12 & 65.59 \\
& ESGPT-30K & 10h 16m 24s & \underline{84.32} & \underline{54.94} & \textbf{73.37} & \underline{65.93} & 59.40 & \underline{63.82} & \underline{66.96} \\
& ESGPT-50K & 17h 26m 58s & \textbf{84.80} & \textbf{55.69} & \underline{71.69} & \textbf{66.30} & 59.01 & \textbf{64.52} & \textbf{67.00} \\
& ESGPT-Full & 146h 29m 35s & 84.23 & 54.60 & 71.19 & 64.20 & 50.15 & 41.73 & 61.02 \\
\hline

\multirow{5}{*}{\makecell[c]{Llama-3.1-\\8B-Instruct}}
& ESGPT-10K & 03h 16m 39s & 81.78 & 49.21 & 69.06 & 61.36 & \underline{55.55} & 58.10 & 62.51 \\
& ESGPT-20K & 07h 20m 15s & 82.25 & 51.72 & 70.35 & 62.35 & \textbf{56.45} & 59.68 & 63.80 \\
& ESGPT-30K & 11h 23m 51s & 82.53 & \underline{52.72} & \textbf{71.97} & \underline{63.83} & 53.23 & \textbf{61.18} & \underline{64.24} \\
& ESGPT-50K & 19h 31m 03s & \textbf{84.14} & \textbf{52.80} & \underline{70.63} & \textbf{67.04} & 53.88 & \underline{59.86} & \textbf{64.72} \\
& ESGPT-Full & 164h 02m 13s & \underline{83.38} & 45.69 & 68.33 & 62.84 & 42.61 & 33.01 & 55.98 \\
\hline
\end{tabular}%
}
\caption{Scale ablation results of \textsf{ESGPT} on \textsf{ESGU-Bench}. The models are grouped by their base model, and each row reports the performance of a training data scale variant. Best scores within each base-model group are shown in bold.}
\label{tab:scale_ablation}
\end{table*}
%------------------------------------
\begin{table*}[t]
\centering
\scriptsize
\setlength{\tabcolsep}{3pt}
\renewcommand{\arraystretch}{1.12}
\begin{tabularx}{\textwidth}{@{}L{0.21\textwidth}L{0.12\textwidth}L{0.07\textwidth}Y L{0.05\textwidth}@{}}
\toprule
\textbf{Dataset} & \textbf{Source Type} & \textbf{Language} & \textbf{Description} & \textbf{Access} \\
\midrule

ESG-FTSE~\cite{pavlova2024esgftse}
& News / Media\newline Articles
& English
& English news dataset on top FTSE 100 firms for ESG relevance detection and primary E/S/G category classification.
& Public \\
\midrule

ESG-Kor~\cite{lee2024esgkor}
& Corporate Reports\newline / Disclosures
& Korean
& Korean ESG sentence dataset from sustainability reports, supporting ESG relevance detection and E/S/G category classification.
& Public \\
\midrule

ESGBERT~\cite{schimanski2024ESGBERT}
& Corporate Reports\newline / Disclosures
& English
& Sentence-level ESG corpus from corporate news, annual reports, and sustainability reports for E/S/G and action classification.
& Public \\
\midrule

ESG Activities~\cite{birti2025optimizing}
& Corporate Reports\newline / Disclosures
& English
& Transportation-sector non-financial disclosure dataset for matching text segments to EU ESG taxonomy activity descriptions.
& Public \\
\midrule

ClimateBERT~\cite{webersinke2021climatebert}
& Climate Reports\newline / Disclosures
& English
& Climate disclosure datasets for climate detection, specificity, sentiment, commitments/actions, TCFD, environmental claims, and target classification.
& Public \\
\midrule

DAX ESG Media~\cite{lipenkova2023detecting}
& News / Media\newline Articles
& English
& DAX company ESG media dataset using filtered external news and web documents for ESG topic extraction.
& Public \\
\midrule

DynamicESG~\cite{Tseng2023dynamicesg}
& News / Media\newline Articles
& Multilingual
& News-based ESG dataset for dynamic ESG category, impact type, and impact duration classification.
& Public \\
\midrule

A3CG~\cite{ong2025towards}
& Corporate Reports\newline / Disclosures
& English
& SGX sustainability report dataset for aspect extraction and action classification of corporate sustainability claims.
& Public \\
\midrule

PromiseEval~\cite{seki2025mlpromise}
& Corporate Reports\newline / Disclosures
& Multilingual
& Multilingual ESG report dataset for promise identification, supporting evidence, evidence clarity, and verification timing.
& Request \\
\midrule

SASB-Aligned ESG Sentences~\cite{junprung2023sasbesg}
& ESG Standards\newline / Frameworks
& English
& Sentence-level ESG dataset aligned with SASB standards for ESG detection and fine-grained SASB issue classification.
& Public \\
\midrule

ESG Sentiment~\cite{trajanovristo2023esgsentiment}
& Corporate Reports\newline / Disclosures
& English
& ESG aspect-level sentiment dataset with E/S/G-specific positive, neutral, and negative labels.
& Public \\
\midrule

ML-ESG-1~\cite{chen2023mlesg}
& News / Media\newline Articles
& Multilingual
& Multilingual ESG news dataset for identifying fine-grained ESG issues based on MSCI ESG key issue guidelines.
& Request \\
\midrule

ML-ESG-2~\cite{chen2023mlesg2}
& News / Media\newline Articles
& Multilingual
& Multilingual ESG dataset for classifying whether ESG-related events imply opportunity, risk, or cannot-distinguish impacts.
& Request \\
\midrule

ML-ESG-3~\cite{chen2024mlesg3}
& News / Media\newline Articles
& Multilingual
& Multilingual ESG news dataset for inferring ESG impact duration and related impact-level/type labels.
& Request \\
\midrule

SustainableQA~\cite{ali2025sustainableqa}
& Corporate Reports\newline / Disclosures
& English
& Large-scale sustainability QA dataset from corporate sustainability and annual reports, covering factoid, non-factoid, and table QA.
& Public \\
\midrule

ESGenius~\cite{he2025esgenius}
& ESG Standards\newline / Frameworks
& English
& Expert-validated ESG knowledge QA benchmark built from major sustainability standards and framework documents.
& Public \\
\midrule

Climate-FEVER~\cite{diggelmann2020climatefever}
& Wikipedia
& English
& Climate claim verification dataset linking real-world climate claims to Wikipedia evidence for support, refute, or NEI labels.
& Public \\

\bottomrule
\end{tabularx}
\caption{ESG resources used for constructing \textsf{ESG-Instruct} and \textsf{ESGU-Bench}.}
\label{tab:esg_resources}
\end{table*}
%------------------------------------
\newcolumntype{L}[1]{>{\RaggedRight\arraybackslash}p{#1}}
\newcolumntype{Y}{>{\RaggedRight\arraybackslash}X}
\newcolumntype{R}[1]{>{\RaggedLeft\arraybackslash}p{#1}}
\newcommand{\NA}{--}

\begin{table*}[t]
\centering
\scriptsize
\setlength{\tabcolsep}{2.5pt}
\renewcommand{\arraystretch}{1.08}
\begin{tabularx}{\textwidth}{
@{}
L{0.18\textwidth}
L{0.13\textwidth}
Y
L{0.065\textwidth}
R{0.065\textwidth}
R{0.065\textwidth}
R{0.065\textwidth}
R{0.065\textwidth}
@{}
}
\toprule
\textbf{Category} & \textbf{Dataset} & \textbf{Task} & \textbf{Language} & \textbf{Train} & \textbf{Valid} & \textbf{Test} & \textbf{Final} \\
\midrule

General
& Alpaca
& General Instruction Chat
& English
& 52000 & \NA & \NA & 3000 \\

\midrule
ESG Relevance Detection
& ESG-Kor
& Korean ESG Relevance Classification
& Korean
& 83262 & 16986 & 16987 & 5160 \\

\midrule
ESG Topic Identification
& A3CG
& Aspect Extraction
& English
& 1673 & 244 & 447 & 575 \\

& DAX ESG Media
& ESG Topic Classification
& English
& 11548 & \NA & \NA & 575 \\

& Dynamic-ESG
& ESG Category Classification
& Chinese
& 1058 & 118 & 131 & 575 \\

& ESG Activities
& ESG Activity--Concept Matching
& English
& 1272 & \NA & 53 & 575 \\

& ESG-BERT Dataset
& Environmental Classification
& English
& 2000 & \NA & \NA & 574 \\

& ESG-BERT Dataset
& Social Classification
& English
& 2000 & \NA & \NA & 574 \\

& ESG-BERT Dataset
& Governance Classification
& English
& 2000 & \NA & \NA & 574 \\

& ESG-Kor
& Korean ESG Category Classification
& Korean
& 79269 & 16986 & 16987 & 574 \\

& ML-ESG1
& ESG Issue Classification
& French
& 1200 & \NA & 300 & 574 \\

\midrule
ESG Commitment Analysis
& A3CG
& Action Classification
& English
& 2336 & 352 & 676 & 2688 \\

& ESG-BERT Dataset
& ESG Action Classification
& English
& 500 & \NA & \NA & 448 \\

& PromiseEval
& Promise Status Classification
& French
& 400 & \NA & 400 & 400 \\

& PromiseEval
& Promise Verification Timeline Classification
& French
& 400 & \NA & 400 & 400 \\

& PromiseEval
& Promise Status Classification
& Japanese
& 400 & \NA & 400 & 400 \\

& PromiseEval
& Promise Verification Timeline Classification
& Japanese
& 400 & \NA & 400 & 400 \\

\midrule
ESG Impact Analysis
& Dynamic-ESG
& Impact Duration Classification
& Chinese
& 995 & 140 & 156 & 470 \\

& Dynamic-ESG
& Impact Type Classification
& Chinese
& 1260 & 111 & 123 & 470 \\

& ESG-Sentiment
& ESG Sentiment Classification
& English
& 611 & \NA & 68 & 470 \\

& ML-ESG2
& Impact Type Classification
& French
& 818 & \NA & 200 & 470 \\

& ML-ESG2
& ESG Sentiment Classification
& Japanese
& 896 & \NA & 225 & 470 \\

& ML-ESG3
& Impact Duration Classification
& French
& 661 & \NA & 146 & 470 \\

& ML-ESG3
& Impact Level Classification
& French
& 661 & \NA & 146 & 470 \\

& ML-ESG3
& Impact Duration Classification
& Japanese
& 1500 & \NA & 53 & 470 \\

& ML-ESG3
& Impact Type Classification
& Japanese
& 1500 & \NA & 53 & 470 \\

& ML-ESG3
& Impact Duration Classification
& Korean
& 800 & \NA & 200 & 470 \\

& ML-ESG3
& Impact Type Classification
& Korean
& 800 & \NA & 200 & 470 \\

\midrule
ESG Claim Verification
& PromiseEval
& Evidence Quality Classification
& French
& 400 & \NA & 400 & 400 \\

& PromiseEval
& Evidence Status Classification
& French
& 400 & \NA & 400 & 400 \\

& PromiseEval
& Evidence Quality Classification
& Japanese
& 400 & \NA & 400 & 400 \\

& PromiseEval
& Evidence Status Classification
& Japanese
& 400 & \NA & 400 & 400 \\

\midrule
ESG Question Answering
& SustainableQA
& Factoid QA
& English
& 71418 & 8835 & 8609 & 1704 \\

& SustainableQA
& Non-Factoid QA
& English
& 82208 & 10140 & 10191 & 2066 \\

& SustainableQA
& Table QA
& English
& 3143 & 392 & 421 & 1394 \\

\bottomrule
\end{tabularx}
\caption{Composition of \textsf{ESG-Instruct}. Train/Valid/Test denote the number of instances in the original source dataset splits when available. Final denotes the number of selected instances used in \textsf{ESG-Instruct} after preprocessing and quality control.}
\label{tab:esg_instruct_composition}
\end{table*}

%------------------------------------
\begin{table*}[t]
\centering
\scriptsize
\setlength{\tabcolsep}{2.5pt}
\renewcommand{\arraystretch}{1.08}
\begin{tabularx}{\textwidth}{
@{}
L{0.18\textwidth}
L{0.13\textwidth}
Y
L{0.065\textwidth}
R{0.065\textwidth}
R{0.065\textwidth}
R{0.065\textwidth}
R{0.065\textwidth}
@{}
}
\toprule
\textbf{Category} & \textbf{Dataset} & \textbf{Task} & \textbf{Language} & \textbf{Train} & \textbf{Valid} & \textbf{Test} & \textbf{Final} \\
\midrule

ESG Relevance Detection
& ESG-FTSE
& ESG Relevance Classification
& English
& 3129 & 390 & 394 & 394 \\

& ClimateBERT
& Climate Detection Classification
& English
& 1300 & \NA & 400 & 400 \\

& ClimateBERT
& Environmental Claims Classification
& English
& 2117 & 265 & 265 & 265 \\

\midrule
ESG Topic Identification
& ESG-FTSE
& ESG Category Classification
& English
& 3127 & 389 & 396 & 396 \\

& ML-ESG1
& ESG Issue Classification
& English
& 1199 & \NA & 300 & 300 \\

& SASB-Aligned ESG Sentences
& SASB Parent Category Classification
& English
& 6460 & \NA & 647 & 647 \\

& SASB-Aligned ESG Sentences
& SASB Issue Topic Classification
& English
& 6460 & \NA & 647 & 647 \\

& ClimateBERT
& TCFD Recommendation Classification
& English
& 1300 & \NA & 400 & 400 \\

\midrule
ESG Commitment Analysis
& ClimateBERT
& Climate Commitments / Actions Classification
& English
& 1000 & \NA & 320 & 320 \\

& ClimateBERT
& Climate Specificity Classification
& English
& 1000 & \NA & 320 & 320 \\

& ClimateBERT
& Net-Zero Reduction Target Classification
& English
& 3441 & \NA & 344 & 344 \\

& PromiseEval
& Promise Status Classification
& English
& 400 & \NA & 400 & 400 \\

& PromiseEval
& Promise Verification Timeline Classification
& English
& 400 & \NA & 400 & 400 \\

\midrule
ESG Impact Analysis
& ClimateBERT
& Climate Sentiment Classification
& English
& 1000 & \NA & 320 & 320 \\

& ML-ESG2
& Impact Type Classification
& English
& 808 & \NA & 218 & 218 \\

& ML-ESG3
& Impact Duration Classification
& English
& 545 & \NA & 136 & 136 \\

& ML-ESG3
& Impact Level Classification
& English
& 545 & \NA & 136 & 136 \\

\midrule
ESG Claim Verification
& Climate Fever
& Claim Verification
& English
& \NA & \NA & 1535 & 1535 \\

& PromiseEval
& Evidence Quality Classification
& English
& \NA & \NA & 400 & 400 \\

& PromiseEval
& Evidence Status Classification
& English
& \NA & \NA & 400 & 400 \\

\midrule
ESG Question Answering
& ESGenius
& ESG Reasoning
& English
& \NA & \NA & 1136 & 1136 \\

\bottomrule
\end{tabularx}
\caption{Composition of \textsf{ESGU-Bench}. Train/Valid/Test denote the number of instances in the original source dataset splits when available. Final denotes the number of selected instances used in the \textsf{ESGU-Bench} evaluation set after preprocessing and quality control.}
\label{tab:esg_benchmark_composition}
\end{table*}

\end{document}

